\documentclass[runningheads,envcountsame]{llncs}

\usepackage[T1]{fontenc}
\usepackage{amsmath}
\usepackage{amssymb}
\let\proof\relax
\let\endproof\relax
\usepackage{amsthm}
\usepackage{microtype}
\usepackage{tikz-cd}
\usepackage{booktabs}
\usepackage{array}
\usepackage[hidelinks,pdfauthor={},pdftitle={Greedy and Matroid Intersection without Branching}]{hyperref}

\newcommand{\R}{\mathbb R}
\newcommand{\Opt}{\operatorname{Opt}}
\newcommand{\cl}{\operatorname{cl}}
\newcommand{\blk}{\operatorname{blk}}
\newcommand{\PO}{\mathrm{PO}}
\newcommand{\Bases}{\mathcal B}
\newcommand{\Pos}{\mathcal P}
\newcommand{\CC}{\mathcal C}
\newcommand{\Av}{\operatorname{Av}}

\spnewtheorem{mainthm}{Theorem}{\bfseries}{\itshape}
\renewcommand{\themainthm}{\Alph{mainthm}}

% Where the Lean development is provided.  If the submission system takes no
% supplementary material, replace by an anonymous link.
\newcommand{\leanartifact}{as supplementary material}

% Line breaking only: a third pass with extra stretch before a line overflows.
\setlength{\emergencystretch}{1em}


\begin{document}

\title{Greedy and Matroid Intersection\texorpdfstring{\\}{ }without Branching}
\titlerunning{Greedy and Matroid Intersection without Branching}
\author{Anonymous Authors}
\authorrunning{Anonymous Authors}
\institute{Anonymous Institution}

\maketitle

\begin{abstract}
The greedy algorithm and the augmenting path algorithms of matroid optimization
compare weights and branch on the outcome.  A $(\max,+,-)$-circuit, like a ReLU
network, computes a continuous function and cannot branch; we show that both
algorithms have exact counterparts of polynomial size.  For a matroid on $m$
elements, a Boolean circuit of size $s$ for the rank function yields a
$(\max,+,-)$-circuit of size $3ms+O(m^{2})$ for the maximum weight of a base, and
the two complexities are polynomially equivalent.  Both are polynomial for binary
matroids, matroids represented over $\mathbb Q$ or a fixed finite field, and
gammoids.  For two matroids with rank circuits of sizes $s_1$ and $s_2$, circuits
of size $O(m^{7}(s_1+s_2+m))$ compute the optimal values of weighted matroid
intersection, including the assignment problem, for which polynomial size was
open.  The circuit updates Lagrange multipliers in a fixed sequence.  The family
of bases these updates should test is NP-hard to query in general, but a
monotone family of matroids containing it keeps the updates valid; a potential
bounds their number.  All circuits and size bounds are verified in Lean~4.
\keywords{Tropical circuits \and ReLU neural networks \and Matroid intersection
\and Greedy algorithm \and Lagrangian relaxation}
\end{abstract}

\section{Introduction}
\label{sec:intro}

The greedy algorithm for weighted bases~\cite{edmonds71} and the augmenting path
algorithms for weighted matroid intersection~\cite{edmonds70,lawler75,frank81}
choose each step by comparing weights.  A circuit over $\max$, $+$ and $-$ cannot
compare: it computes a continuous function of its inputs, while the outcome of a
comparison is discontinuous.  Circuits over $\min$ and $+$ alone need exponential
size for such problems: minimum spanning tree needs
$(\min,+)$-circuits of size $2^{\Omega(\sqrt n)}$ on $n$-vertex
graphs~\cite{jukna-seiwert18-greedy}, and minimum-cost arborescence and the
assignment problem need size $2^{\Omega(n)}$~\cite{jerrum-snir82}.  Jukna surveys these
tropical circuits~\cite{jukna15-tropical,jukna23-tropical}.

With subtraction, Hertrich and Sering~\cite{hertrich-sering25} construct
polynomial-size circuits for the value of a minimum spanning tree and
for maximum flow.  Fomin, Grigoriev and Koshevoy~\cite{fgk16} give small
subtraction-free circuits for the generating polynomials of spanning trees and of
directed spanning trees; tropicalizing them gives polynomial-size
$(\min,+,-)$-circuits for minimum-cost arborescence.  Hertrich,
Kober and Loho~\cite{hertrich-kober-loho26} tropicalize a small circuit for the
basis generating polynomial of a regular matroid and name binary matroids as the
natural next step, noting that counting their bases is \#P-hard and that
different techniques would be needed.  By the identity
$\max(a,b)=\max(0,a-b)+\max(0,b)-\max(0,-b)$, a $(\max,+,-)$-circuit of size $S$
converts into a ReLU neural network of size $3S$~\cite{hertrich-kober-loho26},
so every circuit below also yields an exact ReLU neural network.

Let $M$ be a matroid on a ground set $E$ with $|E|=m$, rank function $r$,
rank $k=r(E)$ and set of bases $\Bases(M)$.  For $u\in\R^{E}$ and $X\subseteq E$ we
write $u(X)=\sum_{e\in X}u_e$; then $h_M(w)=\max_{B\in\Bases(M)}w(B)$ is the largest
weight of a base, and $L=\lceil\log_2(k+1)\rceil$ is the number of binary digits of
$k$.  Our circuits read
$M$ only through a \emph{Boolean rank circuit}, a Boolean circuit that outputs
the binary digits of $r(X)$ from the indicator vector of $X\subseteq E$.

\begin{mainthm}[Maximum-weight bases]
\label{thm:A}
If $M$ has a Boolean rank circuit of size $s$, then a $(\max,+,-)$-circuit of
size at most $3ms+mL+9m^{2}+4m+4L+1$ computes $h_M(w)$ at every $w\in\R^{E}$.
\end{mainthm}

\begin{mainthm}[Converse]
\label{thm:B}
If a $(\max,+,-)$-circuit of size $S$ computes $h_M$ on $\{0,1\}^{E}$, then $M$
has a Boolean rank circuit of size at most $15S^{2}+40S$.
\end{mainthm}

Hence, for every matroid, the $(\max,+,-)$-circuit complexity of $h_M$ and the
Boolean circuit complexity of its rank function are polynomially related.
Binary matroids,
matroids represented over $\mathbb Q$ or over a fixed finite field, graphic
matroids and the duals of all of these (Corollary~\ref{cor:classes}), as well as
transversal matroids and gammoids (Corollary~\ref{cor:transversal} in
Appendix~\ref{app:two}), have Boolean rank circuits of polynomial size, so for
all of them $h_M$ has a $(\max,+,-)$-circuit of polynomial size.  By counting, most
matroids need $(\max,+,-)$-circuits of size $2^{m/2-O(\log m)}$
(Corollary~\ref{cor:count} in Appendix~\ref{app:greedy}).

Let the ground set of $M$ be partitioned into blocks $R_1,\ldots,R_k$.  A set is
\emph{rainbow} if it meets every block at most once; a rainbow base meets every
block exactly once.  Let $\Opt(w)$ be the least weight $w(B)$ of a rainbow base $B$
of $M$.  Rainbow spanning trees are the rainbow
bases of a graphic matroid with colored edges, and their generating polynomial
specializes to the permanent (Proposition~\ref{prop:star} in
Appendix~\ref{app:lagr}).  Arborescences are the rainbow bases of a graphic
matroid whose blocks are the in-arcs of the non-root vertices.

\begin{mainthm}[Rainbow bases]
\label{thm:C}
If $M$ has a Boolean rank circuit of size $s$ and a rainbow base, then a uniform
$(\max,+,-)$-circuit of size $O(k^{3}m^{4}(s+m))$ computes $\Opt(w)$ at every
$w\in\R^{E}$.
\end{mainthm}

\begin{mainthm}[Weighted matroid intersection]
\label{thm:D}
Let $M_1$ and $M_2$ be matroids on $E$, $|E|=m$, with Boolean rank circuits of
sizes $s_1$ and $s_2$.  If they have a common base, a uniform
$(\max,+,-)$-circuit of size $O(m^{7}(s_1+s_2+m))$ computes the minimum weight of
a common base at every $w\in\R^{E}$.  For all such $M_1$ and $M_2$, one of size
$O(m^{7}(s_1+s_2+m))$ computes the maximum weight of a common independent set.
\end{mainthm}

Theorem~\ref{thm:D} follows from Theorem~\ref{thm:C} through the classical
correspondence between common bases of $M_1,M_2$ and rainbow bases of
$M_1\oplus M_2^{*}$ on two copies of $E$ (Section~\ref{sec:two}).

For the $n\times n$ assignment problem, the ground set consists of the $n^{2}$
entries of the cost matrix, $M$ is the partition matroid with one part per
column, which has a Boolean rank circuit of size $O(n^{2})$ (one disjunction per
column and a population count), and the blocks are the rows.  Every connected
component of $M$ is a column of rank one, so Corollary~\ref{cor:rankone} and the
proof of Theorem~\ref{thm:C} give $(\max,+,-)$-circuits, and hence ReLU neural
networks, of size $O(k^{3}m^{3}(s+m))=O(n^{11})$.  Hertrich and
Sering~\cite{hertrich-sering25} and
Jukna~\cite[Section~6.5, Problem~3]{jukna23-tropical} ask whether polynomial
size is possible for this problem, and Hertrich and
Stargalla~\cite{hertrich-stargalla26}, who list it as open, prove that circuits
over $\max$, $+$ and multiplication by positive scalars need $2^{\Omega(n)}$
addition gates for it.

\paragraph{Techniques.}
The circuit of Theorem~\ref{thm:C} maintains potentials $p\in\R^{k}$ on the
blocks and, with $\blk(e)$ the block of $e$, the Lagrangian
\[
  \Psi(p)=\sum_{b}p_b+\min_{B\in\Bases(M)}\sum_{e\in B}\bigl(w_e-p_{\blk(e)}\bigr).
\]
It is at most $\Opt(w)$ and equals it once some minimum base of the shifted
weights $v_e=w_e-p_{\blk(e)}$ is rainbow.  Stage $r=1,\ldots,k-1$ raises the potentials
until a minimum base of $v$ meets $r+1$ blocks.  An update raises $p_b$ by
the length of the set of levels $t\ge0$ at which a matroid $N(t)$ that grows with $t$
has a base that meets exactly $r$ blocks and misses $b$, up to the first level at which $N(t)$
has a base meeting $r+1$ blocks.  This length is an integral over $t$ of
Boolean tests on $N(t)$, which Corollary~\ref{cor:levels} turns into a
$(\max,+,-)$-circuit.

The exact family for these tests is the family $\Pos_t$ of bases that are optimal for
some weight within $t/2$ of $v$ in every coordinate, the possibly optimal bases
of interval data~\cite{kasperski08}.  It is not a matroid, and deciding whether
it contains a rainbow spanning tree is NP-complete
(Proposition~\ref{prop:obstacles}), so we enlarge it.  By
Theorem~\ref{thm:validity}, for every monotone family of matroids that starts at
the minimum bases and contains $\Pos_t$, each update raises $\Psi$ by exactly
$k-r$ times its length, and every minimum base that meets $r$ blocks stays
minimum.  Our family clusters the possibly optimal elements of each connected
component by single linkage at scale $t$ (Proposition~\ref{prop:posl}); it
changes at only $2m$ levels, which the circuit computes from the weights
(Lemma~\ref{lem:events}).  A potential that
counts pairs in which a fully raised element lies at or below a less raised one
bounds the number of updates
(Theorems~\ref{thm:potential} and~\ref{thm:stagelength}).

\paragraph{Related work.}
Hertrich and Loho~\cite{hertrich-loho26} show that the virtual extension
complexity of a polytope is at most the size of any neural network that
optimizes over it, so Theorem~\ref{thm:A} with Corollary~\ref{cor:classes}(1) and
Theorem~\ref{thm:D} bound it by a polynomial for the base polytope of a binary
matroid and for the common base and common independent set polytopes of two
binary matroids.  The circuits of~\cite{fgk16,hertrich-kober-loho26}
tropicalize generating polynomials; ours are built from Boolean rank circuits,
so they cover binary matroids and matroids represented over $\mathbb Q$ or over
a fixed finite field, and extend to weighted matroid intersection.  The maximum flow circuits of~\cite{hertrich-sering25} also run a fixed number
of steps, each of which acts only when a guessed residual distance is correct;
in our circuits the length of each step is an integral of Boolean tests, which
Corollary~\ref{cor:levels} computes.
Edmonds~\cite{edmonds70,edmonds71}, Lawler~\cite{lawler75} and
Frank~\cite{frank81} solve the same problems by comparisons and augmenting
paths.  The correctness and the size bounds of the circuits are formalized in
Lean~4 (Section~\ref{sec:checks}); complete proofs are in the appendix.

\section{Model and Tools}
\label{sec:model}

\begin{definition}
\label{def:circuits}
A \emph{$(\max,+,-)$-circuit} on inputs $x\in\R^{n}$ is a directed acyclic
graph whose sources are the coordinates of $x$ and the constant $0$, and whose
other nodes are gates $\max(a,b)$, $a+b$ and $a-b$.  A \emph{Boolean circuit} on
$x\in\{0,1\}^{n}$ has the coordinates of $x$ and the constant $0$ as sources and
gates $\wedge$, $\vee$ and $\neg$.  In both, the size is the number of gates.  A
Boolean circuit \emph{outputs the binary digits} of an integer if its output
bits, least significant first, are these digits.  A family of circuits
is \emph{uniform} if a polynomial-time algorithm produces each member from its
combinatorial data, here the ground set, the blocks and the given Boolean rank
circuits.
\end{definition}

We use standard matroid terminology~\cite{oxley11,schrijver03} and write
$\cl$ for the closure operator of $M$, $X+e=X\cup\{e\}$, $X-e=X\setminus\{e\}$, $[k]=\{1,\ldots,k\}$, $\mathbf 1_X$ for the indicator vector of
a set $X$, $\mathbf 1=\mathbf 1_E$ and $\mathbf 1[P]\in\{0,1\}$ for the truth value of a
statement $P$.  For a base $B$
and $e\notin B$, $C(B,e)$ is the fundamental circuit in $B+e$.  For $u\in\R^{E}$
write $\beta_M(u)=\min_Bu(B)$.  The minimum bases of
$u$ are the bases of a matroid $M^{u}$, and they decompose over the connected
components of $M$ (Lemma~\ref{lem:minbases} in Appendix~\ref{app:model}).

\begin{definition}
\label{def:ti}
For $D\in\R^{E}_{\ge0}$ and $t\ge0$ let $G_t(D)=\{e:D_e\le t\}$, and let
$0=d_0\le d_1\le\cdots\le d_m$ be the entries of $D$ in increasing order with a
zero prepended.  Let $\Delta_i=d_{i+1}-d_i$.  For
$F:2^{E}\to\R$ let
\[
  \CC_F(D)=\sum_{i=0}^{m-1}\Delta_i\,F\bigl(G_{d_i}(D)\bigr)
  =\int_0^{d_m}F\bigl(G_t(D)\bigr)\,dt.
\]
\end{definition}

If $F(E)=0$, then $\CC_F(D)=\int_0^{\infty}F(G_t(D))\,dt$ is the Lov\'asz
extension of $X\mapsto F(E\setminus X)$ at $D$~\cite{lovasz83}.

\begin{theorem}
\label{thm:compiler}
Let a single Boolean circuit of size $s$ compute $F=(F_1,\ldots,F_N):2^{E}\to\{0,1\}^{N}$
from the indicator vector of its argument.  Then a $(\max,+,-)$-circuit of
size at most $3ms+mN+6m^{2}$ computes the values $\CC_{F_j}(D)$,
$j=1,\ldots,N$, at every $D\in\R^{E}_{\ge0}$.
\end{theorem}

\begin{proof}[Proof sketch]
A sorting network gives the levels $d_i$ and the differences $\Delta_i$.  Since $D_e$ occurs among
the $d_i$, we have $\max(\Delta_i-\max(D_e-d_i,0),0)=\Delta_i\mathbf 1[e\in G_{d_i}(D)]$,
and for $x,y\in\{0,1\}$ and $\Delta\ge0$,
\[
  \max(\Delta x+\Delta y-\Delta,0)=\Delta(x\wedge y),\quad
  \max(\Delta x,\Delta y)=\Delta(x\vee y),\quad
  \Delta-\Delta x=\Delta(\neg x).
\]
So the Boolean circuit, simulated gate by gate, outputs $\Delta_iF_j(G_{d_i}(D))$,
and the sums over $i$ are $\CC_{F_j}(D)$ (Appendix~\ref{app:model}).
\end{proof}

Taking the levels and the tested values from different vectors gives the
following corollary.

\begin{corollary}
\label{cor:levels}
Let $\Lambda\in\R^{n}_{\ge0}$, let $0=d_0\le d_1\le\cdots\le d_n$ be its entries in
increasing order with a zero prepended and $\Delta_i=d_{i+1}-d_i$, and let
$Y\in\R^{q}$ have every entry among $d_0,\ldots,d_n$.  Let $\mathbf 1[Y\le d_i]$ have the entries $\mathbf 1[Y_l\le d_i]$.  If a single
Boolean circuit of size $s$ computes $F=(F_1,\ldots,F_N):\{0,1\}^{q}\to\{0,1\}^{N}$,
then a $(\max,+,-)$-circuit of size at most $\frac32n^{2}+n(1+4q+3s+N)$ computes
from $\Lambda$ and $Y$ the values
\[
  \sum_{i=0}^{n-1}\Delta_i\,F_j\bigl(\mathbf 1[Y\le d_i]\bigr),\qquad j=1,\ldots,N.
\]
\end{corollary}

\section{Greedy without Branching}
\label{sec:greedy}

The greedy algorithm reads the weights in decreasing order; its value is the
Lov\'asz extension of the rank function~\cite{edmonds70,lovasz83}, which has the
form of Definition~\ref{def:ti}.

\begin{lemma}
\label{lem:greedy}
For every $D\in\R^{E}_{\ge0}$ and $\mu\in\R$,
\begin{align*}
  h_M(D)&=\sum_{i=0}^{m-1}\Delta_i\,r\bigl(E\setminus G_{d_i}(D)\bigr)=\CC_F(D)
  \quad\text{with }F(X)=r(E\setminus X),\\
  h_M(D+\mu\mathbf 1)&=h_M(D)+\mu\,r(E).
\end{align*}
\end{lemma}

\begin{proof}[Proof sketch]
With $X_i=E\setminus G_{d_i}(D)$, so that $D_e=\sum_{i:\,e\in X_i}\Delta_i$, every
base $B$ satisfies
\[
  D(B)=\sum_{i=0}^{m-1}\Delta_i\,|B\cap X_i|\le\sum_{i=0}^{m-1}\Delta_i\,r(X_i),\qquad
  (D+\mu\mathbf 1)(B)=D(B)+\mu\,r(E).
\]
A base that extends nested bases of $X_{m-1}\subseteq\cdots\subseteq X_0$
attains the bound.
\end{proof}

\begin{proof}[Proof sketch of Theorems~\ref{thm:A} and~\ref{thm:B}]
For Theorem~\ref{thm:A}, Lemma~\ref{lem:greedy} with $\mu=\min_ew_e$ and
$F(X)=r(E\setminus X)$ gives $h_M(w)=\mu k+\CC_F(w-\mu\mathbf 1)$, and
Corollary~\ref{cor:integer} in Appendix~\ref{app:model} computes $\CC_F$.  For
Theorem~\ref{thm:B}, on $\{0,1\}^{E}$ the wires of a $(\max,+,-)$-circuit with $S$
gates carry integers in $[-2^{S},2^{S}]$, so $(S+2)$-bit two's complement
arithmetic simulates it gate by gate, and $h_M(\mathbf 1_T)=r(T)$
(Appendix~\ref{app:greedy}).
\end{proof}

Gaussian elimination without branching, and for duals the formula
$r^{*}(X)=|X|+r(E\setminus X)-r(E)$, give the rank circuits of the next corollary.

\begin{corollary}
\label{cor:classes}
For the following matroids $M$ on $m$ elements, the function $h_M$ has a $(\max,+,-)$-circuit of polynomial size.
\begin{enumerate}
  \item A binary matroid represented by an $n\times m$ matrix over
        $\mathrm{GF}(2)$: size at most $100(m+1)^{2}(n+1)^{2}$.
  \item A matroid represented by an $n\times m$ matrix over a finite field with
        $q$ elements: size at most $25q^{2}(m+1)^{2}(n+1)^{2}$.
  \item A matroid represented over $\mathbb Q$ by an $n\times m$ integer matrix
        whose entries have absolute value less than $2^{b}$: size at most
        $2^{21}(m+1)^{2}(n+1)^{2}(n(n+b)+1)^{6}$.
  \item The dual $M^{*}$ of a matroid $M$ with a Boolean rank circuit of size
        $s$: size $3ms+O(m^{2}\log m)$.
  \item The cycle matroid $M(G)$ of a multigraph $G$ with $n$ vertices and $m$
        edges, and its dual: sizes at most $100(m+1)^{2}(n+1)^{2}$ and
        $100(m+1)^{2}((n+1)^{2}+m+1)$.
\end{enumerate}
\end{corollary}

\section{Rainbow Bases}
\label{sec:rainbow}

\subsection{Setting and the Lagrangian}
\label{sec:lagr}

Throughout this section we write $r_M$ for the rank function of $M$, as the letter
$r$ denotes a stage, and $E_Q=\bigcup_{b\in Q}R_b$ for $Q\subseteq[k]$.  For a
matroid $N$ on $E$ let $\nu(N)$ be the
largest size of a rainbow independent set of $N$, and let $\Av(N)$ be the set of
blocks missed by some rainbow independent set of size $\nu(N)$, the
\emph{avoidable} blocks.  We assume that $M$ has a rainbow base; the algorithm
that produces the circuit checks this once by unweighted matroid intersection.  For
$y\in\R^{k}$ let $y_{\blk}\in\R^{E}$ have the entries $y_{\blk(e)}$.  For
$p\in\R^{k}$ the shifted weights are $v=w-p_{\blk}$ and the Lagrangian is
$\Psi(p)=\sum_{b=1}^{k}p_b+\beta_M(v)$.  Let
\[
  \nu_0=\nu(M^{v}),\qquad S_0=\Av(M^{v}).
\]
So $\nu_0$ is the largest number of blocks met by a minimum base of $v$.

\begin{lemma}
\label{lem:duality}
$\Psi(p)\le\Opt(w)$ for every $p$, and $\Psi(p)=\Opt(w)$ whenever $\nu_0=k$.
\end{lemma}

\subsection{Possibly Optimal Bases and Two Obstacles}
\label{sec:pos}

\begin{definition}
\label{def:theta}
For a base $B$ let
$\theta(B)=\max\bigl(\{0\}\cup\{v_f-v_e:e\notin B,\ f\in C(B,e)\}\bigr)$, the
largest gain of a single exchange, and let
$\Pos_t=\{B\in\Bases(M):\theta(B)\le t\}$ for $t\ge0$.
\end{definition}

\begin{proposition}
\label{prop:pos}
For a base $B$ and $t\ge0$ the following are equivalent.
\begin{enumerate}
  \item $\theta(B)\le t$.
  \item $B$ is a minimum base for $v-t\mathbf 1_B$.
  \item $B$ is a minimum base for some $u$ with $\|u-v\|_\infty\le t/2$.
\end{enumerate}
\end{proposition}

So $\Pos_t$ is the family of bases that are possibly optimal for the interval
weights $[v-\tfrac t2,v+\tfrac t2]$~\cite{kasperski08}.  It grows with $t$,
equals $\Bases(M^{v})$ at $t=0$, and equals $\Bases(M)$ for
$t\ge\max v-\min v$.

\begin{proposition}
\label{prop:obstacles}
\begin{enumerate}
  \item For $U_{2,4}$ on $\{a,b,c,d\}$ with $v=(0,1,2,3)$, the family $\Pos_2$
        equals $\{ab,ac,ad,bc\}$ and violates the basis exchange axiom.
  \item Given a multigraph with colored edges, edge weights $v$ and $t\ge0$,
        deciding whether $\Pos_t$ contains a rainbow spanning tree is
        NP-complete.
\end{enumerate}
\end{proposition}

\subsection{Admissible Families and the Update}
\label{sec:enlarge}

We therefore use matroids $N(t)$ with $\Pos_t\subseteq\Bases(N(t))$.

\begin{definition}
\label{def:admissible}
A family of matroids $N(t)$ on $E$, $t\ge0$, is \emph{admissible} if
(F1)~$N(0)=M^{v}$; (F2)~every base of $N(t)$ is a base of $N(t')$ whenever
$t\le t'$; (F3)~$\Pos_t\subseteq\Bases(N(t))$ for every $t$; (F4)~$N(t)$ has a
rainbow base for all large $t$; and (F5)~there are levels
$0=c_0<c_1<\cdots<c_q$ such that $N(t)$ is constant on each $[c_i,c_{i+1})$ and
on $[c_q,\infty)$.
\end{definition}

\begin{definition}
\label{def:update}
Let $r<k$ and $\nu_0\ge r$.  For an admissible family let $\tau$ be the least
$t$ with $\nu(N(t))\ge r+1$, let $\lambda_b$ be the least $t$ at which some base
of $N(t)$ meeting exactly $r$ blocks misses $b$, with $\lambda_b=\infty$ if there
is none, and let $a_b=(\tau-\lambda_b)_+=\max(\tau-\lambda_b,0)$.  The vector $a$ is
the update of stage $r$, and stage $r$ replaces $p$ by $p+a$ while $\nu_0=r$.  If $\nu_0>r$, then $\tau=0$ and $a=0$.
\end{definition}

\begin{lemma}
\label{lem:updint}
Let $r<k$, $\nu_0\ge r$, and let $N(t)$ be admissible.  The minima in
Definition~\ref{def:update} exist, $\nu(N(t))=r$ and
$\Av(N(t))=\{b:\lambda_b\le t\}$ for $0\le t<\tau$, and
\[
  a_b=(\tau-\lambda_b)_+
  =\int_0^{\infty}\mathbf 1\bigl[\nu(N(t))\le r\bigr]\;\mathbf 1\bigl[b\in\Av(N(t))\bigr]\,dt.
\]
\end{lemma}

\begin{theorem}[Validity]
\label{thm:validity}
Let $\nu_0=r<k$, let $N(t)$ be admissible, and let $a$ be its update.  Then
\begin{enumerate}
  \item $\Psi(p+a)=\Psi(p)+(k-r)\tau$;
  \item every base of $M^{v}$ meeting $r$ blocks is a minimum base of
        $v-a_{\blk}$; in particular $\nu_0$ does not decrease.
\end{enumerate}
\end{theorem}

\begin{proof}[Proof sketch]
For $t\in[0,\tau]$ let $a(t)_b=(t-\lambda_b)_+$, $v(t)=v-a(t)_{\blk}$ and
$S_t=\{b:\lambda_b\le t\}$.  The function $\varphi(t)=\Psi(p+a(t))$ is piecewise
linear with right derivative
$|S_t|-\max\{|B\cap E_{S_t}|:B\in\Bases(M^{v(t)})\}$.  Let $t<\tau$.  As
$a(t)\in[0,t]^{k}$, every minimum base of $v(t)$ lies in $\Pos_t$
(Lemma~\ref{lem:path} in Appendix~\ref{app:pos}), hence in $\Bases(N(t))$ by
(F3).  By Lemma~\ref{lem:updint}, $\nu(N(t))=r$ and $\Av(N(t))=S_t$, and by
Lemma~\ref{lem:avoidable} in Appendix~\ref{app:update},
$|\Av(N)|-r_N(E_{\Av(N)})=k-\nu(N)$ for every matroid $N$.  A base $B^{*}$ of
$M^{v}$ meeting $r$ blocks misses only blocks of $S_0\subseteq S_t$, so
\[
  \max_{B\in\Bases(M^{v(t)})}|B\cap E_{S_t}|\le r_{N(t)}(E_{S_t})=|S_t|-(k-r)
  \le|B^{*}\cap E_{S_t}|.
\]
Hence $\varphi'(t^{+})\ge k-r$, while $\varphi^{*}(t)=\sum_b(p_b+a(t)_b)+v(t)(B^{*})$
satisfies $\varphi^{*}\ge\varphi$, $\varphi^{*}(0)=\varphi(0)$ and
$(\varphi^{*})'(t^{+})\le k-r$ (Appendix~\ref{app:update}).  So $B^{*}$ stays
minimum and
\[
  \Psi(p)+(k-r)t\le\varphi(t)\le\varphi^{*}(t)\le\Psi(p)+(k-r)t.\qedhere
\]
\end{proof}

\subsection{Single Linkage on Possibly Optimal Elements}
\label{sec:posl}

\begin{definition}
\label{def:posl}
For $t\ge0$ let $\PO_t=\bigl\{f\in E:f\notin\cl\{g:v_g<v_f-t\}\bigr\}$.  In each
connected component $K$ of $M$, link $x,y\in\PO_t\cap K$ when $|v_x-v_y|\le t$,
and call the classes of the transitive closure of this relation the
\emph{clusters} of $K$ at scale $t$.  Each cluster is an interval of the sorted
values of $\PO_t\cap K$ in which consecutive values differ by at most $t$, and
consecutive values in different clusters differ by more than $t$, so the clusters
of $K$ are ordered.  For $e\in\PO_t$ let $u_t(e)$ be the position of the cluster
of $e$ in this order, $1$ for the lowest cluster of its component, let
$u_t(e)=m+1$ for $e\notin\PO_t$, and let $N(t)=M^{u_t}$.  We call $N$ the
\emph{PO-SL} family, for single linkage on possibly optimal elements.
\end{definition}

By Lemma~\ref{lem:minbases}(4), $N(t)$ depends only on the comparisons of $u_t$
within components.  Write $\mathcal G_t=\{(e,f)\in E\times E:v_f-v_e\le t\}$.

\begin{proposition}
\label{prop:posl}
For every $v$, the family $N(t)$ of Definition~\ref{def:posl} satisfies
(F1)--(F5) with levels $\{0\}\cup\{|v_e-v_f|:e,f\in E\}$, and it is determined
by $\mathcal G_t$.
\end{proposition}

\begin{proof}[Proof sketch]
For (F3), a base with $\theta(B)\le t$ is minimum for $v-t\mathbf 1_B$, so
$B\subseteq\PO_t$, and an exchange gain $v_f-v_e\in(0,t]$ links $e$ and $f$.  For
(F2), let $t<t'$ and let $B$ be a minimum base of $u_t$.  As the scale grows,
$\PO_t$ grows and clusters merge in order, so it remains to consider
$e\in\PO_{t'}\setminus\PO_t$ outside $B$ and $f\in C(B,e)$ with $v_f>v_e$.  Then
$e\in\cl\{h\in B:u_t(h)\le u_t(x)\}$ for some $x\in\PO_t$ with $v_x<v_e-t$, and as
clusters are ordered by value (Appendix~\ref{app:posl}),
\[
  u_t(f)\le u_t(x),\ v_x<v_e<v_f\ \Rightarrow\ u_t(f)=u_t(x)\ \Rightarrow\
  u_{t'}(f)=u_{t'}(x)=u_{t'}(e).\qedhere
\]
\end{proof}

\subsection{Breakpoints}
\label{sec:events}

By Proposition~\ref{prop:posl}, $N(t)$ is determined by the $m^{2}$ bits of
$\mathcal G_t$, which change at the $m^{2}$ levels $|v_e-v_f|$;
Lemma~\ref{lem:events} reduces them to $2m$ levels.  Call two elements,
not necessarily distinct, \emph{joined} if they are non-loops in the same connected component of $M$, and
fix a linear order on $E$.

\begin{lemma}
\label{lem:events}
Fix $v\in\R^{E}$.
\begin{enumerate}
  \item For a non-loop $x$ let $D^{x}_g=\max(v_x-v_g,0)$ and
        $\phi_x(S)=\mathbf 1[x\in\cl(E\setminus S)]$, and put
        $t_x=\CC_{\phi_x}(D^{x})$; put $t_x=0$ for a loop.  For $t\ge0$, a non-loop
        $x$ lies in $\PO_t$ if and only if $t\ge t_x$, and a loop lies in no
        $\PO_t$.
  \item For joined $x,y$ let $W_{xx}=t_x$ and $W_{xy}=\max(t_x,t_y,|v_x-v_y|)$ for
        $x\ne y$, and let $T_{xy}$ be the least value of
        $\max_iW_{a_ia_{i+1}}$ over the sequences $x=a_0,a_1,\ldots,a_n=y$ of joined
        elements with $n\ge1$.  For $t\ge0$, $x$ and $y$ lie in the same cluster at scale
        $t$ if and only if $T_{xy}\le t$.
  \item Let $\sigma_x$ be the least $T_{xy}$ over the $y<x$ joined to $x$, or
        $0$ if there is none, and let $\Lambda$ consist of the $2m$ values
        $t_x$ and $\sigma_x$, $x\in E$.  Then $T_{xy}\in\Lambda\cup\{0\}$ for all
        joined $x,y$.
  \item For joined $x,y$ let
        $\omega_{xy}=\max\bigl(t_x,t_y,\min(T_{xy},v_y-v_x)\bigr)$.  Then
        $\omega_{xy}\in\Lambda\cup\{0\}$, and for $t\ge0$, $u_t(x)<u_t(y)$ if and
        only if $t_x\le t<\omega_{xy}$.
\end{enumerate}
In particular $N(t)$ is constant on every interval between consecutive values of
$\Lambda\cup\{0\}$ and beyond the largest one.
\end{lemma}

\begin{proof}[Proof sketch]
An element $x$ enters $\PO_t$ at $t_x=\CC_{\phi_x}(D^{x})$, where $\phi_x$ is
antitone, and two clusters merge at the minimax distance $T_{xy}$.
If no element of the cluster of $x$ at scale $h=T_{xy}$ enters at $h$, then
$T_{ef}<h$ is an equivalence relation on this cluster, and the least element $z$ of
a class that avoids the least element of the cluster has $\sigma_z=h$
(Appendix~\ref{app:events}).
\end{proof}

\subsection{Termination}
\label{sec:term}

Fix a stage $r<k$ and a state, that is, potentials $p$ with $\nu_0=r$.  Let $a$ be
the update of PO-SL at this state, and let $v'=v-a_{\blk}$.  The blocks of $S_0$ are raised by the full
length $\tau$ of the update, and the others by less.  Let $\Pi(v)$ be the set of pairs in which an element of a fully raised block
lies at or below an element of another block,
\[
  \Pi(v)=\bigl\{(y,g):y\in E_{S_0},\ g\notin E_{S_0},\ v_y\le v_g\bigr\}.
\]

\begin{theorem}[Potential]
\label{thm:potential}
The following hold.
\begin{enumerate}
  \item If $\nu(M^{v'})=r$, then $\Av(M^{v'})\supseteq S_0$.
  \item $\Pi(v')\supseteq\Pi(v)$, both formed with $S_0$.
  \item $\Pi(v')\ne\Pi(v)$.
\end{enumerate}
Consequently stage $r$ ends after at most $r\bigl(\lfloor m^{2}/4\rfloor+1\bigr)$
updates.
\end{theorem}

\begin{proof}[Proof sketch]
Item 1 follows from Theorem~\ref{thm:validity}(2), and item 2 from
$a_b=\tau\ge a_{b'}$ for $b\in S_0$ and every block $b'$.  For item 3, let
$\delta>0$ be the first level at which $\Av(N(t))$ or $\nu(N(t))$ changes, so
that $a_b\le\tau-\delta$ for $b\notin S_0$.  By Lemma~\ref{lem:avoidable},
$r_{N(\delta)}(E_{S_0})>|S_0|-(k-r)=r_{M^{v}}(E_{S_0})$, so for a minimum base $B$
of $v$ with $|B\cap E_{S_0}|=r_{M^{v}}(E_{S_0})$,
some $g\in B\setminus E_{S_0}$ and $y\in E_{S_0}\setminus B$ lie in the same cluster
at scale $\delta$ with $v_g<v_y$.  In that cluster, some consecutive values
$\alpha=v_{z_1}<\beta=v_{z_2}$ with
$z_1\notin E_{S_0}$ and $z_2\in E_{S_0}$ differ by at most $\delta$, and
\[
  v'_{z_2}-v'_{z_1}\le(\beta-\tau)-\bigl(\alpha-(\tau-\delta)\bigr)=(\beta-\alpha)-\delta\le0,
  \qquad(z_2,z_1)\in\Pi(v')\setminus\Pi(v).
\]
Within a stage, $S_0$ only grows and $k-r\le|S_0|\le k-1$, so it takes at most
$r$ values, and $|\Pi|\le|E_{S_0}|\,|E\setminus E_{S_0}|\le\lfloor m^{2}/4\rfloor$
(Appendix~\ref{app:term}).
\end{proof}

\begin{theorem}[Stage length]
\label{thm:stagelength}
Stage $r$ ends after at most $\min\bigl(1+(r-1)km,\,1+k(2m-k-1)\bigr)$ updates.
\end{theorem}

\begin{proof}[Proof sketch]
Fix a base $B$ of $M^{v}$ meeting $r$ blocks at the start of the stage; by
Theorem~\ref{thm:validity} it remains one.  By Lemma~\ref{lem:crossing}(1) in Appendix~\ref{app:term}, each
update moves some $x\in E_{S_0}$ to or below some $z\notin E_{S_0}$ with $x$ or
$z$ in $B$, so every update but the last enlarges
the set of pairs $(x,z)$, $x\ne z$, with $x$ or $z$ in $B$, $\blk(x)\in\Av(M^{v})$
and $v_x\le v_z$, which has at most $k(m-1)+(m-k)k=k(2m-k-1)$ elements; counting
within the runs of constant $S_0$ gives $1+(r-1)km$ (Appendix~\ref{app:term}).
\end{proof}

\begin{corollary}
\label{cor:rankone}
If every connected component of $M$ has rank at most one, every update of stage $r$
changes $S_0$ or ends the stage, and stage $r$ ends after at most $r$ updates.
\end{corollary}

\begin{proof}[Proof sketch]
For a base $B$ of $M^{v}$ meeting $r$ blocks, Lemma~\ref{lem:crossing}(1) in
Appendix~\ref{app:term} gives $x\in E_{S_0}$ and $z\notin E_{S_0}$ in the same
component, one of them in $B$, with $v_x>v_z$ and $v'_x\le v'_z$.  So $\{x,z\}$
is a circuit, and as $B$ is minimum for $v$ and for $v'$
(Theorem~\ref{thm:validity}), $z\in B$, $x\notin B$ and
\[
  v'_z\le v'_x\le v'_z,\qquad |(B-z+x)\cap E_{S_0}|=|S_0|-(k-r)+1.
\]
By Lemma~\ref{lem:avoidable}, this minimum base $B-z+x$ of $v'$ makes $\nu_0$
rise or $S_0$ grow (Theorem~\ref{thm:potential}(1)), and $k-r\le|S_0|\le k-1$.
\end{proof}

\subsection{The Circuit}
\label{sec:circuit}

The tests behind an update are questions of unweighted matroid
intersection, which the augmenting path algorithm answers without branching on
weights.

\begin{lemma}
\label{lem:augment}
Let $N_x$ be a matroid on $E$ for every vector $x$ of input bits, and for each
$y\in E$ let a Boolean circuit of size $q$ decide, from $x$ and the indicator
vector of a set $J$ with $y\in J$ and $J-y$ independent in $N_x$, whether $J$ is
independent in $N_x$.  For $j\ge0$, a Boolean circuit of size
$O(j(m^{2}q+m^{3}))$ computes from $x$ the indicator vector of a rainbow
independent set of $N_x$ with at least $\min(j,\nu(N_x))$ elements, and with
$\nu(N_x)$ elements for $j=k$.  For $r_0\ge0$, a Boolean circuit of size
$O((r_0+1)m^{2}(q+m))$ computes from $x$ the $k$ bits
$\mathbf 1[\nu(N_x)\le r_0]\wedge\mathbf 1[b\in\Av(N_x)]$.
\end{lemma}

\begin{lemma}
\label{lem:observer}
Let $M$ have a Boolean rank circuit of size $s$, and fix
$r_0\in\{1,\ldots,k-1\}$.  For
$v\in\R^{E}$ let $N(t)$ be the family of Definition~\ref{def:posl} and
$\mathcal Q(v)$ the vector of the $m+2m^{2}$ values $t_x$,
$T_{xy}$ and $\omega_{xy}$ of Lemma~\ref{lem:events}, with $T_{xy}=\omega_{xy}=0$ for
$x,y$ not joined.  A Boolean circuit of size $O(r_0m^{2}(s+m))$ reads $m+2m^{2}$
bits and outputs $k$ bits $F_b$, $b\in[k]$, such that for every $v\in\R^{E}$ and
$t\ge0$, on input $\mathbf 1[\mathcal Q(v)\le t]$,
\[
  F_b=\mathbf 1\bigl[\nu(N(t))\le r_0\bigr]\wedge\mathbf 1\bigl[b\in\Av(N(t))\bigr].
\]
\end{lemma}

\begin{proof}[Proof sketch]
By Lemma~\ref{lem:events}, the bits $\mathbf 1[\mathcal Q(v)\le t]$ give $\PO_t$
and the clusters at no cost, and their order with $2m^{2}$ gates.  Let
$\mathrm{Lev}(y)$ be the cluster of $y$ and $\mathrm{Below}(y)$ the set of
elements of its component with smaller $u_t$.  By the rank formula of $M^{u_t}$,
for a set $A$ with $A-y$ independent in $N(t)$,
\[
  \begin{aligned}
  &A\text{ independent in }N(t)\iff y\in\PO_t\ \text{and}\\
  &\qquad r_M\bigl(\mathrm{Below}(y)\cup(A\cap\mathrm{Lev}(y))\bigr)
  =r_M\bigl(\mathrm{Below}(y)\bigr)+|A\cap\mathrm{Lev}(y)|.
  \end{aligned}
\]
So each test costs one rank evaluation, and Lemma~\ref{lem:augment} gives the
bits $F_b$ (Appendix~\ref{app:circuit}).
\end{proof}

\begin{lemma}
\label{lem:oneupdate}
Let $M$ have a Boolean rank circuit of size $s$.  For each stage $r$, a
$(\max,+,-)$-circuit of size $O(rm^{3}(s+m))$ maps
$(w,p)$ to $p+a$, with $a$ the PO-SL update at $v=w-p_{\blk}$, at every state
with $\nu_0\ge r$.
\end{lemma}

\begin{proof}[Proof sketch]
The circuit computes $v$, the entry levels $t_x$ by
Theorem~\ref{thm:compiler}, the merge levels by the Floyd--Warshall algorithm
over $(\min,\max)$, whose bits $\mathbf 1[T_{xy}\le t]$ follow Warshall's
algorithm~\cite{warshall62} at every $t$, and from them $\Lambda$ and
$\mathcal Q(v)$.  With $0=d_0\le\cdots\le d_{2m}$ the entries of $\Lambda$ in
increasing order with a zero prepended and $F_b$
the circuit of Lemma~\ref{lem:observer} for $r_0=r$, Lemmas~\ref{lem:updint}
and~\ref{lem:events} give
\[
  a_b=\sum_{i=0}^{2m-1}(d_{i+1}-d_i)\,F_b\bigl(\mathbf 1[\mathcal Q(v)\le d_i]\bigr).
\]
Corollary~\ref{cor:levels} computes this sum with $2m$ simulated copies of that
circuit (Appendix~\ref{app:circuit}).
\end{proof}

\begin{proof}[Proof of Theorem~\ref{thm:C}]
Start from $p=0$, where $\nu_0\ge1$, and for $r=1,\ldots,k-1$ apply the circuit of
Lemma~\ref{lem:oneupdate} for stage $r$, which is the identity at $\nu_0>r$,
exactly $\ell_r$ times, where
\[
  \ell_r=\min\bigl(r(\lfloor m^{2}/4\rfloor+1),\,1+(r-1)km,\,1+k(2m-k-1)\bigr).
\]  As $\nu_0$ never decreases
(Theorem~\ref{thm:validity}) and Theorems~\ref{thm:potential}
and~\ref{thm:stagelength} bound the number of updates at $\nu_0=r$ by $\ell_r$,
we have $\nu_0=k$ at the end.  By Lemma~\ref{lem:duality} and
Theorem~\ref{thm:A}, $O(ms+m^{2})$ further gates then output
\[
  \Psi(p)=\sum_{b=1}^{k}p_b+\beta_M(v)=\sum_{b=1}^{k}p_b-h_M(-v)=\Opt(w).
\]
An update of stage $r$ has $O(rm^{3}(s+m))$ gates, and $\ell_r\le2km$ gives
\[
  \sum_{r=1}^{k-1}r\,\ell_r\le2km\sum_{r=1}^{k-1}r=k^{2}(k-1)m\le k^{3}m.
\]
So the circuit has size $O(k^{3}m^{4}(s+m))$.  If every connected component of $M$ has
rank at most one, Corollary~\ref{cor:rankone} allows $\ell_r=r$, and
$\sum_rr\,\ell_r\le k^{3}$ gives the size $O(k^{3}m^{3}(s+m))$.  Uniformity is
shown in Appendix~\ref{app:circuit}.
\end{proof}

\section{Two Matroids}
\label{sec:two}

\begin{lemma}
\label{lem:doubling}
Let $M_1,M_2$ be matroids of rank $\rho$ on $E$.  On $E\times\{1,2\}$ let
$N=M_1\oplus M_2^{*}$, with blocks $\{(e,1),(e,2)\}$ and weight $w_e$ on $(e,1)$,
$0$ on $(e,2)$.  For $X\subseteq E$, the set $X\times\{1\}\cup(E\setminus X)\times\{2\}$,
which meets every block once, is a base of $N$ if and only if $X$ is a common base
of $M_1$ and $M_2$, and its weight is $w(X)$.  Every rainbow base of $N$ has this
form.
\end{lemma}

\begin{proof}[Proof sketch of Theorem~\ref{thm:D}]
The matroid $N$ of Lemma~\ref{lem:doubling} has rank $m$ on $2m$ elements with
$m$ blocks, and its rank function costs $s_1+s_2+O(m)$ gates since, with $r_i$
the rank function of $M_i$,
\[
  r_N\bigl(X_1\times\{1\}\cup X_2\times\{2\}\bigr)=r_1(X_1)+|X_2|+r_2(E\setminus X_2)-r_2(E).
\]
So Theorem~\ref{thm:C} with $k=m$ gives the minimum
weight of a common base with $O(m^{3}(2m)^{4}(s_1+s_2+m))=O(m^{7}(s_1+s_2+m))$
gates.  For common independent sets $X$, the truncations $M_1',M_2'$ at $m$ of the
direct sums of $M_1$ and $M_2$ with a free matroid on $m$ new elements $Z$ have the
common base $Z$, and with $w'$ equal to $-w$ on $E$ and $0$ on $Z$,
$\max_Xw(X)=-\min\{w'(B):B\in\Bases(M_1')\cap\Bases(M_2')\}$
(Appendix~\ref{app:two}).
\end{proof}

\section{Formal Verification}
\label{sec:checks}

A Lean~4 development with Mathlib, both at version 4.32.0, proves the results of
Sections~\ref{sec:model} to~\ref{sec:two}, with explicit bounds in place of
$O(\cdot)$, for the matroids of Mathlib's \texttt{Matroid} on $\{0,\ldots,m-1\}$;
there a circuit is a list of gates with the semantics of
Definition~\ref{def:circuits}, and its size is the number of gates.  Uniformity is
not formalized, and of Proposition~\ref{prop:obstacles}(2) only the reduction is.  The
main theorems depend on no axioms beyond the
three standard axioms of Lean.  The development is provided \leanartifact;
Table~\ref{tab:lean} in Appendix~\ref{app:checks} lists the results it proves.

\section{Open Problems}
\label{sec:scope}

By Theorem~\ref{thm:stagelength}, a stage of PO-SL takes $O(km)$ updates.  Does
it take $O(m)$, or can it take $\Theta(km)$?  Is there an admissible family determined by
$\mathcal G_t$ whose stages take $O(r)$ updates?  For the assignment
problem, how far can the exponent $11$ be lowered?  Do rainbow questions about
$\Pos_t$ stay NP-hard on the states that these updates reach, and is there a
min-max formula of Rado type for $\Pos_t$ on such states?  Can a fixed sequence of such updates
compute weighted matching in general graphs, an instance of matroid parity whose
linear programming dual has a variable for every odd set of vertices?  Are there optimization problems solvable in
polynomial time whose values need $(\max,+,-)$-circuits of super-polynomial
size?

\bibliographystyle{splncs04}
\bibliography{gi_ipco}

\clearpage
\appendix
% Generated by build_appendix.py from appendix_src.tex and
% ../greedy_intersection.tex.  Edit those files, not this one.

\section{Proofs for Section~\texorpdfstring{\ref{sec:model}}{2}}
\label{app:model}

We use the following classical facts~\cite{oxley11,schrijver03}.

\begin{lemma}
\label{lem:minbases}
Let $u\in\R^{E}$.
\begin{enumerate}
  \item A base $B$ is minimum for $u$ if and only if $u_f\le u_e$ whenever
        $e\notin B$ and $f\in C(B,e)$.
  \item A base $B$ is minimum for $u$ if and only if $B\cap\{u\le c\}$ spans
        $\{u\le c\}$ for every $c\in\R$.
  \item The minimum bases of $u$ are the bases of a matroid $M^{u}$.  For a
        minimum base $B$, $e\notin B$ and $f\in B$, the set $B-f+e$ is a minimum
        base if and only if $f\in C(B,e)$ and $u_f=u_e$, so the circuit of $e$ in
        $B+e$ in $M^{u}$ is $\{e\}\cup\{f\in C(B,e):u_f=u_e\}$.
  \item The relation ``$e=f$, or $e$ and $f$ lie on a common circuit'' is an
        equivalence relation, whose classes are the \emph{connected components};
        a loop or a coloop is a component by itself.  Closure, rank and minimum
        bases decompose over components: $e\in\cl(X)$ if and only if
        $e\in\cl(X\cap K)$ for the component $K$ of $e$, $r(X)=\sum_Kr(X\cap K)$,
        and a base $B$ is minimum if and only if every $B\cap K$ is a minimum base
        of $M|K$.  In particular the minimum bases of $u$, and $M^{u}$, depend
        only on the comparisons of $u$ within components.
\end{enumerate}
\end{lemma}

Explicitly, $M^{u}$ is the direct sum, over the values $\ell$ of $u$, of the
matroids $(M|\{u\le\ell\})/\{u<\ell\}$ on the level sets $\{u=\ell\}$, so its
rank function is read from that of $M$: if $\ell_1<\cdots<\ell_q$ are the values
of $u$ and $U_i=\{u\le\ell_i\}$, $U_0=\varnothing$, then
\begin{equation}
  r_{M^{u}}(A)=\sum_{i=1}^{q}\Bigl(r\bigl(U_{i-1}\cup(A\cap\{u=\ell_i\})\bigr)-r(U_{i-1})\Bigr).
  \label{eq:rank-mu}
\end{equation}

The two expressions in Definition~\ref{def:ti} agree, as $G_t(D)=G_{d_i}(D)$ for
$t\in[d_i,d_{i+1})$.  If $F(E)=0$, the upper limit $d_m=\max_eD_e$ can be replaced
by $\infty$, as $G_t(D)=E$ for $t\ge d_m$.

\begin{proof}[Proof of Theorem~\ref{thm:compiler}]
\emph{Step 1: $d_i$ and $\Delta_i$.}  The insertion sorting network~\cite{knuth98-taocp3} has
$m(m-1)/2$ comparators, each computing
\[
  \max(a,b)\quad\text{and}\quad\min(a,b)=a+b-\max(a,b)
\]
with three gates.  It outputs $d_1,\ldots,d_m$, and one gate per $i<m$ gives
$\Delta_i=d_{i+1}-d_i$.

\emph{Step 2: $\Delta_i\mathbf 1[e\in G_{d_i}(D)]$.}  For $i<m$ and $e\in E$, the entry $D_e$
occurs in the sorted list, so $D_e\le d_i$ or $D_e\ge d_{i+1}$, and
\begin{equation}
  \max\bigl(\Delta_i-\max(D_e-d_i,0),\,0\bigr)
  =\begin{cases}\Delta_i,&D_e\le d_i\\0,&D_e\ge d_{i+1}\end{cases}
  =\Delta_i\,\mathbf 1\bigl[e\in G_{d_i}(D)\bigr],
  \label{eq:bit}
\end{equation}
at four gates.

\emph{Step 3: $\Delta_iF_j(G_{d_i}(D))$.}  For $\Delta=\Delta_i\ge0$ and $x,y\in\{0,1\}$,
\[
  \max(\Delta x+\Delta y-\Delta,0)=\Delta(x\wedge y),\quad
  \max(\Delta x,\Delta y)=\Delta(x\vee y),\quad
  \Delta-\Delta x=\Delta(\neg x).
\]
Let $B$ be the given Boolean circuit and $B_\Delta$ the circuit in which each gate of
$B$ is replaced by the left side, at most three gates.  By induction on the gates,
\[
  \begin{tikzcd}[column sep=4.5em,row sep=2.2em]
    \{0,1\}^{E}\arrow[r,"B"]\arrow[d,"\Delta\,\cdot"'] & \{0,1\}^{N}\arrow[d,"\Delta\,\cdot"]\\
    \R^{E}\arrow[r,"B_\Delta"'] & \R^{N}
  \end{tikzcd}
\]
commutes, and at the inputs of \eqref{eq:bit}
\[
  B_\Delta\bigl(\Delta\mathbf 1_{G_{d_i}(D)}\bigr)=\Delta B\bigl(\mathbf 1_{G_{d_i}(D)}\bigr)
  =\Delta_iF\bigl(G_{d_i}(D)\bigr).
\]

\emph{Step 4: $\CC_{F_j}(D)$ with at most $3ms+mN+6m^{2}$ gates.}  With $N$ additions for each $i$,
\[
  \sum_{i=0}^{m-1}\Delta_iF_j\bigl(G_{d_i}(D)\bigr)=\CC_{F_j}(D),\qquad j=1,\ldots,N,
\]
and the size is at most
\[
  \underbrace{\tfrac32m(m-1)}_{\text{sorting}}
  +m\bigl(\underbrace{1}_{\Delta_i}+\underbrace{4m}_{\eqref{eq:bit}}
  +\underbrace{3s}_{\text{simulation}}+\underbrace{N}_{\text{sums}}\bigr)
  \le3ms+mN+6m^{2}.\qedhere
\]
\end{proof}

\begin{corollary}
\label{cor:integer}
Let $f:2^{E}\to\mathbb N$, and let a Boolean circuit of size $s$ output the
binary digits of $f$ on $L$ wires.  Then a $(\max,+,-)$-circuit of size at
most $3ms+mL+6m^{2}+2L$ computes $\CC_f(D)$ at every $D\in\R^{E}_{\ge0}$.
\end{corollary}

\begin{proof}
\emph{Step 1: $\CC_{f_j}(D)$ for the digits $f_j$ of $f$.}  Write $f=\sum_{j=0}^{L-1}2^{j}f_j$ with
$f_j:2^{E}\to\{0,1\}$ the outputs of the given circuit.  Theorem~\ref{thm:compiler}
with $N=L$ gives $c_j=\CC_{f_j}(D)$, $0\le j<L$, at $3ms+mL+6m^{2}$ gates.

\emph{Step 2: $\CC_f(D)=\sum_{j<L}2^{j}\CC_{f_j}(D)$.}  Put $y_L=0$ and $y_j=c_j+(y_{j+1}+y_{j+1})$, two
gates per digit.  Since $F\mapsto\CC_F(D)$ is linear,
\[
  y_0=\sum_{j=0}^{L-1}2^{j}c_j=\CC_{\sum_j2^{j}f_j}(D)=\CC_f(D).\qedhere
\]
\end{proof}

\begin{proof}[Proof of Corollary~\ref{cor:levels}]
Step 1 of the proof of Theorem~\ref{thm:compiler} sorts $\Lambda$.  In Step 2, for
$i<n$ and an entry $Y_l$, which occurs among $d_0,\ldots,d_n$, either $Y_l\le d_i$
or $Y_l\ge d_{i+1}$, so \eqref{eq:bit} holds with $Y_l$ in place of $D_e$:
\[
  \max\bigl(\Delta_i-\max(Y_l-d_i,0),\,0\bigr)=\Delta_i\,\mathbf 1[Y_l\le d_i].
\]
Steps 3 and 4 are unchanged, with $n$ differences $\Delta_i$ and $q$ indicator values for each $i$, and
the count is
\[
  \tfrac32n(n-1)+n\bigl(1+4q+3s+N\bigr)\le\tfrac32n^{2}+n\bigl(1+4q+3s+N\bigr).\qedhere
\]
\end{proof}

\section{Proofs for Section~\texorpdfstring{\ref{sec:greedy}}{3}}
\label{app:greedy}

\begin{proof}[Proof of Lemma~\ref{lem:greedy}]
\emph{Step 1: $D(B)=\sum_i\Delta_i\,|B\cap X_i|$.}  Let $X_i=E\setminus G_{d_i}(D)=\{e:D_e>d_i\}$, and for
$e\in E$ let $d_{i_e}$ be the first occurrence of $D_e$ in the sorted list.  Then
$e\in X_i$ if and only if $i<i_e$, so
\[
  D_e=d_{i_e}-d_0=\sum_{i<i_e}\Delta_i=\sum_{i:\,e\in X_i}\Delta_i,
  \quad
  D(B)=\sum_{e\in B}\ \sum_{i:\,e\in X_i}\Delta_i=\sum_{i=0}^{m-1}\Delta_i\,|B\cap X_i|
\]
for every base $B$.

\emph{Step 2: $D(B)\le\sum_i\Delta_i\,r(X_i)$.}  Since $B\cap X_i$ is independent and $\Delta_i\ge0$,
\[
  D(B)=\sum_{i=0}^{m-1}\Delta_i\,|B\cap X_i|\le\sum_{i=0}^{m-1}\Delta_i\,r(X_i).
\]

\emph{Step 3: $h_M(D)=\sum_i\Delta_i\,r(X_i)=\CC_F(D)$.}  As $X_0\supseteq\cdots\supseteq X_{m-1}$,
there are bases $I_i$ of $X_i$ with $I_{m-1}\subseteq\cdots\subseteq I_0$; extend
$I_0$ to a base $B$.  Then $I_i\subseteq B\cap X_i$, and
\[
  r(X_i)=|I_i|\le|B\cap X_i|\le r(X_i),
\]
so $B$ attains the bound of Step 2, and by Definition~\ref{def:ti}
\[
  h_M(D)=D(B)=\sum_{i=0}^{m-1}\Delta_i\,r(X_i)=\CC_F(D).
\]

\emph{Step 4: $h_M(D+\mu\mathbf 1)=h_M(D)+\mu\,r(E)$.}  Every base has $r(E)$ elements, so
\[
  h_M(D+\mu\mathbf 1)=\max_{B}\bigl(D(B)+\mu\,|B|\bigr)=h_M(D)+\mu\,r(E).\qedhere
\]
\end{proof}

\begin{proof}[Proof of Theorem~\ref{thm:A}]
\emph{Step 1: $D=w-\mu\mathbf 1\ge0$ with $4m$ gates.}  Put $\mu=\min_ew_e$ and $D=w-\mu\mathbf 1\ge0$, at
$3(m-1)+m\le4m$ gates.

\emph{Step 2: the digits of $F(X)=r(E\setminus X)$ with $s+m$ gates.}  Let $F(X)=r(E\setminus X)$.  The given circuit
on the negated input bits has size at most $s+m$, and since $F(X)\le k<2^{L}$,
its first $L$ outputs are the binary digits of $F(X)$.

\emph{Step 3: $\CC_F(D)$ with $3m(s+m)+mL+6m^{2}+2L$ gates.}  This is Corollary~\ref{cor:integer} for the circuit of Step 2, with $s+m$ in
place of $s$.

\emph{Step 4: $\mu k+\CC_F(D)=h_M(w)$.}  Doubling along the binary digits of $k$ gives $\mu k$
at $2L$ gates, and one addition gives, by Lemma~\ref{lem:greedy},
\[
  \mu k+\CC_F(D)=h_M(D)+\mu\,r(E)=h_M(D+\mu\mathbf 1)=h_M(w).
\]

\emph{Step 5: at most $3ms+mL+9m^{2}+4m+4L+1$ gates.}  The size is at most
\[
  4m+\bigl(3m(s+m)+mL+6m^{2}+2L\bigr)+(2L+1)=3ms+mL+9m^{2}+4m+4L+1.\qedhere
\]
\end{proof}

\begin{proof}[Proof of Theorem~\ref{thm:B}]
\emph{Step 1: $|a_j|\le2^{S}<2^{W-1}$.}  On inputs in $\{0,1\}$ the value $a_j$ of the $j$-th gate
is an integer.  Its arguments $a,b$ are sources or earlier gates, so
$|a|,|b|\le2^{j-1}$ by induction, and as $|a\pm b|\le2\max(|a|,|b|)$ and
$|\max(a,b)|\le\max(|a|,|b|)$,
\[
  |a_j|\le2\max(|a|,|b|)\le2^{j}\le2^{S}<2^{W-1}\qquad(W=S+2).
\]
So all values are exact in $W$-bit two's complement, where sums and differences
are exact and two's complement comparison is the comparison of integers.

\emph{Step 2: each gate with at most $15W+10$ Boolean gates.}  With full adders of eleven gates and $1=\neg0$, the
gadgets cost at most
\[
  a+b:\ 11W,\qquad a-b=a+\neg b+1:\ 12W+1,\qquad \max(a,b):\ 15W+10,
\]
where, with $c$ the carry out of $a+\neg b+1$,
\[
  \neg[a<b]=(a_{W-1}\oplus b_{W-1})\oplus c,\qquad
  \max(a,b)_i=\bigl([a<b]\wedge b_i\bigr)\vee\bigl(\neg[a<b]\wedge a_i\bigr).
\]

\emph{Step 3: $15S^{2}+40S$ gates in all.}  Each input bit becomes a $W$-bit number with high bits
$0$, and each of the $S$ gates its gadget, at most
\[
  S(15W+10)=S\bigl(15(S+2)+10\bigr)=15S^{2}+40S
\]
gates in all; by Step 1 the wires carry the values of the $(\max,+,-)$-circuit.

\emph{Step 4: $h_M(\mathbf 1_T)=r(T)$.}  For $T\subseteq E$,
\[
  h_M(\mathbf 1_T)=\max_{B}|B\cap T|=r(T),\qquad0\le r(T)\le2^{S}<2^{W-1},
\]
so the output bits are the binary digits of $r(T)$.
\end{proof}

\begin{proof}[Proof of Corollary~\ref{cor:classes}]
\emph{Item 1.}  The circuit keeps $n$ slots, slot $p$ holding an occupancy bit
$f_p$ and a vector $b_p\in\mathrm{GF}(2)^{n}$, all empty at the start, and
processes column $e$ with its input bit $x_e$ by
\begin{align*}
  &v_i\leftarrow x_e\wedge A_{ie}\qquad(i\le n),\\
  &v_i\leftarrow v_i\oplus\bigl((v_p\wedge f_p)\wedge b_{p,i}\bigr)\qquad(p=1,\ldots,n,\ i\le n),\\
  &\ell_p=v_p\wedge\neg\textstyle\bigvee_{p'<p}v_{p'},\qquad f_p\leftarrow f_p\vee\ell_p,\qquad
   b_{p,i}\leftarrow(\ell_p\wedge v_i)\vee(\neg\ell_p\wedge b_{p,i}).
\end{align*}
Let $P=\{p:f_p=1\}$.  By induction on the columns, after column $e$ every $b_p$ with
$p\in P$ has $b_{p,i}=0$ for $i<p$ and $b_{p,p}=1$, and
\[
  \operatorname{span}\{b_p:p\in P\}=\operatorname{span}\{A_{\cdot e'}:e'\le e,\ x_{e'}=1\},
\]
since after the second line
\[
  v\in x_eA_{\cdot e}+\operatorname{span}\{b_p:p\in P\},\qquad v_p=0\quad(p\in P),
\]
and the third line stores $v$, if $v\ne0$, in the slot $\min\{i:v_i=1\}\notin P$.
Vectors with leading $1$s in distinct positions are independent, so at the end
the number of occupied slots, $|P|=\dim\operatorname{span}\{A_{\cdot e}:x_e=1\}$, is
the rank of the input set, and the size is
\[
  m(9n^{2}+5n)+11n\lceil\log_2(n+1)\rceil,
\]
the last term for counting the occupied slots.  Theorem~\ref{thm:A} with $k\le n$ gives the bound of item 1.

\emph{Item 2.}  The same elimination runs over $\mathrm{GF}(q)$, an element
being carried by $q$ wires of which exactly one is true: a binary field operation
costs $2q^{2}$ gates, inversion none, masking by a bit $q$ and selection $3q$.
The column step
\[
  \begin{gathered}
  v\leftarrow x_eA_{\cdot e},\qquad v\leftarrow v-[f_p]\,v_p\,b_p\ \ (p=1,\ldots,n),\\
  f_p\leftarrow f_p\vee\ell_p,\qquad
  b_p\leftarrow\begin{cases}v_p^{-1}v&\text{if }\ell_p,\\b_p&\text{otherwise}\end{cases}
  \end{gathered}
\]
costs $6q^{2}n^{2}+3qn^{2}+n^{2}+qn+6n+1$ gates and keeps the invariant of item 1
with leading entries $1$.  Counting costs $11n\lceil\log_2(n+1)\rceil$ gates and the constant $1$ one gate, so the rank circuit has
\[
  m(6q^{2}n^{2}+3qn^{2}+n^{2}+qn+6n+1)+11n\lceil\log_2(n+1)\rceil+1
\]
gates, and Theorem~\ref{thm:A}, applied to this count with $k\le n$ and $q\ge2$, gives the bound of item 2.

\emph{Item 3.}  Put $d=2n(n+b)+1$ and $N=2^{2\lceil\log_2(d+2)\rceil+1}$, so that
$N\le8(d+1)^{2}$ and $(N+1)N^{d}<2^{N}$.  For $S\subseteq E$, with $A_S$ the
matrix of the columns in $S$, and a prime $p$,
$\operatorname{rk}_{\mathrm{GF}(p)}A_S\le\operatorname{rk}_{\mathbb Q}A_S$, since a
rational dependence scales to an integer one with coprime coefficients, which
stays nontrivial modulo $p$.  Conversely let $I$ be a basis of $S$ over
$\mathbb Q$ and $g=\det(A_I^{\mathsf T}A_I)$.  By Hadamard's inequality,
\[
  0<|g|\le\prod_{j\in I}\|A_{\cdot j}\|_2^{2}\le(n\,4^{b})^{n}\le2^{2n(n+b)}<2^{d}.
\]
If every prime $p\le N$ divided $g$, then $2^{\pi(N)}\le\prod_{p\le N}p\le|g|<2^{d}$,
and since $\binom Nj$ divides
$\operatorname{lcm}(1,\ldots,N)=\prod_{p\le N}p^{\lfloor\log_pN\rfloor}$,
\[
  2^{N}=\sum_{j=0}^{N}\binom Nj\le(N+1)\operatorname{lcm}(1,\ldots,N)
  \le(N+1)N^{\pi(N)}<(N+1)N^{d}<2^{N}.
\]
So some prime $p\le N$ does not divide $g$, $A_I$ stays independent modulo $p$, and
\[
  \operatorname{rk}_{\mathbb Q}A_S=|I|\le\operatorname{rk}_{\mathrm{GF}(p)}A_S
  \le\max_{p'\le N\text{ prime}}\operatorname{rk}_{\mathrm{GF}(p')}A_S\le\operatorname{rk}_{\mathbb Q}A_S,
\]
so the maximum is $\operatorname{rk}_{\mathbb Q}A_S$.
The circuits of item 2 for all primes $p\le N$ with a running maximum, by the
two's complement comparison of the proof of Theorem~\ref{thm:B}, have at most
$2^{19}(m+1)(n+1)^{2}(n(n+b)+1)^{6}$ gates, and Theorem~\ref{thm:A} with $k\le n$ gives the bound of item 3.

\emph{Item 4.}  With $W=\lceil\log_2(m+1)\rceil$,
\[
  \begin{aligned}
  r^{*}(S)&=|S|+r(E\setminus S)-r(E)\\
  &\equiv\bigl(2^{W}-(r(E)\bmod2^{W})\bigr)+\sum_{e\in S}1+r(E\setminus S)\pmod{2^{W}},
  \end{aligned}
\]
computed by a $W$-bit counter started at the constant and fed the $m$ input bits
($11Wm$ gates), a ripple-carry adder for $r(E\setminus S)$ ($11W$ gates), $m$
negations, the given circuit and one gate for the constant $1$.  Since
$0\le r^{*}(S)\le m<2^{W}$, the result is $r^{*}(S)$, at $s+m+1+11(m+1)W$ gates, and Theorem~\ref{thm:A} for $M^{*}$ gives the bound of item 4.

\emph{Item 5.}  The vertex-edge incidence matrix of $G$, an $n\times m$ matrix over
$\mathrm{GF}(2)$, represents $M(G)$~\cite{oxley11}.  Item 1 gives the first bound, and the rank circuits of the proofs of items 1 and 4 with Theorem~\ref{thm:A} give the second.
\end{proof}

\begin{corollary}
\label{cor:count}
For all large $m$, most matroids $M$ on $m$ labeled elements have no $(\max,+,-)$-circuit for $h_M$ of size less than $2^{m/2}/m^{2}$.
\end{corollary}

\begin{proof}
\emph{Step 1: at least $2^{2^{m}/(m(m+1))}$ matroids on $E$.}  Let $m\ge2$, $k=\lfloor m/2\rfloor$,
$E=\{0,\ldots,m-1\}$, and let $\mathcal H$ be a largest class of $k$-subsets of
$E$ under $X\mapsto\sum_{x\in X}x\bmod m$, so
$|\mathcal H|\ge\binom mk/m$~\cite{graham-sloane80}.  Distinct $X,Y\in\mathcal H$
satisfy $|X\cap Y|\le k-2$: otherwise $X\setminus Y=\{a\}$, $Y\setminus X=\{b\}$
and
\[
  a-b\equiv\sum_{x\in X}x-\sum_{y\in Y}y\equiv0\pmod m,
\]
so $a=b$.  For $H\subseteq\mathcal H$ let $\mathcal B_H=\binom Ek\setminus H$.  It
is nonempty: for $X\in H$, $a\in X$ and $b\notin X$, $|(X-a+b)\cap X|=k-1$, so
$X-a+b\notin H$.  It satisfies basis exchange: for
$B_1,B_2\in\mathcal B_H$ and $e\in B_1\setminus B_2$, either
$B_2\setminus B_1=\{f\}$ and $B_1-e+f=B_2$, or $B_2\setminus B_1\ni f_1\ne f_2$ and
\[
  \bigl|(B_1-e+f_1)\cap(B_1-e+f_2)\bigr|=k-1,
\]
so $B_1-e+f_1$ or $B_1-e+f_2$ lies in $\mathcal B_H$.  Distinct $H$ give distinct
matroids, and $\binom mk\ge2^{m}/(m+1)$ as the largest of $m+1$ coefficients with
sum $2^{m}$, so the number of matroids on $E$ is at least
\[
  2^{|\mathcal H|}\ge2^{\binom mk/m}\ge2^{2^{m}/(m(m+1))}.
\]
Knuth~\cite{knuth74} proves the stronger bound $2^{2^{m-\frac32\log_2m-O(1)}}$.

\emph{Step 2: at most $2^{O(2^{m}m^{-5/2})}$ of them have small circuits.}  Boolean circuits of size $s'$ on $m$ inputs
with at most $m$ outputs number $2^{O(s'\log(s'+m))}$, and each computes at most
one rank function.  If $h_M$ has a $(\max,+,-)$-circuit of size
$S<2^{m/2}/m^{2}$, Theorem~\ref{thm:B}, keeping the first $m$ output bits, gives a
Boolean rank circuit of size
\[
  s'\le15S^{2}+40S<\frac{15\cdot2^{m}}{m^{4}}+\frac{40\cdot2^{m/2}}{m^{2}}<\frac{2^{m}}{m^{7/2}}
\]
for large $m$, and
\[
  2^{O(s'\log(s'+m))}=2^{O(2^{m}m^{-5/2})}=o\bigl(2^{2^{m}/(m(m+1))}\bigr).\qedhere
\]
\end{proof}

\section{Proofs for Section~\texorpdfstring{\ref{sec:rainbow}}{4}}
\label{app:rainbow}

\subsection{Setting and the Lagrangian}
\label{app:lagr}

The definitions of Section~\ref{sec:lagr} and Lemmas~\ref{lem:avoidable}
and~\ref{lem:augment} apply verbatim to every partition of a finite ground set
into any number $k$ of blocks; Appendix~\ref{app:two} uses them in this form.
Rainbow independent sets of $M^{v}$ extend to minimum bases of $v$, and a minimum
base meeting $j$ blocks contains a rainbow set of size $j$; so $\nu_0$ is the
largest number of blocks met by a minimum base of $v$, and $S_0$ is the set of
blocks missed by some minimum base meeting $\nu_0$ blocks.

\begin{proof}[Proof of Lemma~\ref{lem:duality}]
\emph{Step 1: $\Psi(p)\le\Opt(w)$.}  A rainbow base $B$ meets every block once, so
\[
  w(B)=\sum_{e\in B}\bigl(v_e+p_{\blk(e)}\bigr)=\sum_{b=1}^{k}p_b+v(B)
  \ge\sum_{b=1}^{k}p_b+\beta_M(v)=\Psi(p).
\]

\emph{Step 2: $\Psi(p)=\Opt(w)$ if $\nu_0=k$.}  If $\nu_0=k$, a rainbow independent set of
$M^{v}$ of size $k$ is a rainbow base $B$ with $v(B)=\beta_M(v)$, and
\[
  \Opt(w)\le w(B)=\sum_{b=1}^{k}p_b+\beta_M(v)=\Psi(p)\le\Opt(w).\qedhere
\]
\end{proof}

\begin{proposition}
\label{prop:star}
For every bipartite graph $H=(U,V,F)$ with $|U|=|V|=n$ there is a multigraph with
colored edges whose rainbow spanning trees are the perfect matchings of $H$, with
the same weights.  Its rainbow spanning tree polynomial is the permanent of the
biadjacency matrix of $H$ with variable entries.
\end{proposition}

\begin{proof}
\emph{Step 1: the multigraph on $\{z\}\cup V$.}  On the vertex set $\{z\}\cup V$ take, for each edge
$ij\in F$ with $i\in U$ and $j\in V$, an edge $zj$ of color $i$ and weight
$w_{ij}$.

\emph{Step 2: $T_\pi$ is rainbow $\iff$ $\{\pi(j)j:j\in V\}$ is a perfect matching.}  Each $j\in V$ is adjacent only to $z$, so the
spanning trees are the sets $T_\pi$ containing, for each $j\in V$, the edge $zj$ of
color $\pi(j)$, where $\pi:V\to U$ and $\pi(j)j\in F$.  As $|U|=|V|$,
\[
  T_\pi\text{ rainbow}\iff\pi\text{ bijective}
  \iff\{\pi(j)j:j\in V\}\text{ is a perfect matching of }H,
\]
and $w(T_\pi)=\sum_{j\in V}w_{\pi(j)j}$ is the weight of this matching.  With
variable weights $x_{ij}$,
\[
  \sum_{T_\pi\text{ rainbow}}\ \prod_{j\in V}x_{\pi(j)j}
  =\sum_{\pi:V\to U\text{ bijective}}\ \prod_{j\in V}X_{\pi(j)j}=\operatorname{per}(X),
\]
where $X$ is the biadjacency matrix with $X_{ij}=x_{ij}$ on the edges and
$X_{ij}=0$ elsewhere, so that a bijection using a non-edge contributes $0$.
\end{proof}

\subsection{Possibly Optimal Bases and Two Obstacles}
\label{app:pos}

\begin{proof}[Proof of Proposition~\ref{prop:pos}]
Throughout, $e\notin B$ and $f\in C(B,e)-e\subseteq B$, and $u=v-t\mathbf 1_B$.

\emph{(1)$\Leftrightarrow$(2).}  By Lemma~\ref{lem:minbases}(1),
\[
  \text{(2)}\iff u_f-u_e=v_f-v_e-t\le0\ \text{for all }e\notin B,\ f\in C(B,e)-e\iff\theta(B)\le t.
\]

\emph{(2)$\Rightarrow$(3).}  The weight $u+\tfrac t2\mathbf 1$ has the minimum
bases of $u$, and
$\|u+\tfrac t2\mathbf 1-v\|_\infty=\|\tfrac t2\mathbf 1-t\mathbf 1_B\|_\infty=\tfrac t2$.

\emph{(3)$\Rightarrow$(1).}  If $B$ is minimum for $u'$ with
$\|u'-v\|_\infty\le t/2$, then $u'_f\le u'_e$ by Lemma~\ref{lem:minbases}(1), and
\[
  v_f-v_e\le\bigl(u'_f+\tfrac t2\bigr)-\bigl(u'_e-\tfrac t2\bigr)\le t.\qedhere
\]
\end{proof}

\begin{lemma}
\label{lem:path}
If $y\in[0,t]^{k}$, every minimum base of $v-y_{\blk}$ lies in $\Pos_t$.
\end{lemma}

\begin{proof}
Let $B$ be a minimum base of $v-y_{\blk}$, $e\notin B$ and $f\in C(B,e)$.  By
Lemma~\ref{lem:minbases}(1), $v_f-y_{\blk(f)}\le v_e-y_{\blk(e)}$, so
\[
  v_f-v_e\le y_{\blk(f)}-y_{\blk(e)}\le t,
\]
and $\theta(B)\le t$.
\end{proof}

\begin{proof}[Proof of Proposition~\ref{prop:obstacles}(1)]
\emph{Step 1: $\Pos_2=\{ab,ac,ad,bc\}$.}  In $U_{2,4}$, $C(B,e)=B+e$ for
$e\notin B$, so
\[
  \begin{gathered}
  \theta(B)=\max\Bigl(0,\ \max_{f\in B}v_f-\min_{e\notin B}v_e\Bigr),\\
  \bigl(\theta(ab),\theta(ac),\theta(ad),\theta(bc),\theta(bd),\theta(cd)\bigr)=(0,1,2,2,3,3).
  \end{gathered}
\]

\emph{Step 2: $ad-a+b,\ ad-a+c\notin\Pos_2$.}  For $B_1=ad$, $B_2=bc$ and
$a\in B_1\setminus B_2$, both candidates $B_1-a+f$, $f\in B_2\setminus B_1$, have
$\theta(bd)=\theta(cd)=3>2$.
\end{proof}

\begin{proof}[Proof of Proposition~\ref{prop:obstacles}(2)]
\emph{Step 1: the problem is in NP.}  For a rainbow spanning tree $T$,
$\theta(T)=\max\bigl(\{0\}\cup\{v_g-v_e:e\notin T,\ g\in C(T,e)\}\bigr)$ is
computed from the fundamental cycles in polynomial time.

\emph{Step 2: a reduction from rainbow path.}  Rainbow path is NP-complete~\cite{cfmy09}: given a graph $H$ with colored edges and vertices $a,b$, is there a
path from $a$ to $b$ with distinct colors?  Let $Q$ be the colors of $H$,
$q=|Q|$, and $h\ge q$ the number of vertices of $H$ (add isolated vertices if
needed).  On $V(H)\cup\{z,y\}$ take the edges of $H$ with weight $5$, an edge
$e^{*}=ab$ of weight $0$ and a new color $c^{*}$, an edge $f=ya$ of weight $0$ and
color $c^{*}$, and for every $x\in V(H)$ and every color in $Q\cup X$ an edge $zx$
of that color and weight $10$, where $X$ is a set of $h-q$ new colors.  There are $h+2$ vertices and $h+1$ colors, so a
rainbow spanning tree uses every color once.  Take $t=5$.

\emph{Step 3: a rainbow $T$ with $\theta(T)\le5$ gives a rainbow path in $H$.}  Let $T$ be a rainbow spanning tree with
$\theta(T)\le5$.  As $f$ is the only edge at $y$ and has the color of $e^{*}$,
$f\in T$ and $e^{*}\notin T$.  Every edge $g$ of the path $P_T$ of $T$ from $a$
to $b$ lies in $C(T,e^{*})$, so
\[
  v_g=v_g-v_{e^{*}}\le\theta(T)\le5<10.
\]
So $P_T$ avoids $z$ and the leaf $y$, lies in $H$, and is rainbow since $T$ is.

\emph{Step 4: a rainbow path in $H$ gives a rainbow $T$ with $\theta(T)\le5$.}  Let $P$ be a rainbow path from $a$ to $b$
in $H$.  Let $T$ consist of $E(P)$, $f$, an edge $zx$ for each
$x\in V(H)\setminus V(P)$ and one edge $zx$ with $x\in V(P)$, these
$h-|E(P)|$ edges at $z$ taking the $h-|E(P)|$ colors of $Q\cup X$ not on $P$.
Then $T$ is a rainbow spanning tree with $C(T,e^{*})=P+e^{*}$ and edge weights at
most $10$, so
\[
  \theta(T)\le\max\bigl(\underbrace{5-0}_{e=e^{*}},\ \underbrace{10-5}_{e\in E(H)\setminus T},\
  \underbrace{10-10}_{e=zx\notin T}\bigr)=5.\qedhere
\]
\end{proof}

\subsection{The Update}
\label{app:update}

\begin{proof}[Proof of Lemma~\ref{lem:updint}]
By (F2) independent sets of $N(t)$ stay independent in $N(t')$ for $t\le t'$, and
$\nu(N(0))=\nu_0\ge r$ by (F1), so
\[
  r\le\nu(N(0))\le\nu(N(t))\le\nu(N(t'))\qquad(t\le t').
\]
Moreover $\nu(N(t))=k>r$ for large $t$ by (F4), and by (F5) both sets of levels in
Definition~\ref{def:update} are unions of intervals $[c_i,c_{i+1})$ and
$[c_q,\infty)$, so the minima exist.  For $t<\tau$ we have $\nu(N(t))=r$; a base of
$N(t)$ meeting exactly $r$ blocks contains a rainbow independent set of size $r$
missing the same blocks, and a rainbow independent set of size $r$ extends to a
base of $N(t)$ meeting no further block.  With (F2) this gives, for $t<\tau$,
\[
  \begin{aligned}
  b\in\Av(N(t))&\iff\text{some base of }N(t)\text{ meeting exactly }r\text{ blocks misses }b\\
  &\iff\lambda_b\le t,
  \end{aligned}
\]
and as $\mathbf 1[\nu(N(t))\le r]=\mathbf 1[t<\tau]$,
\[
  \int_0^{\infty}\mathbf 1\bigl[\nu(N(t))\le r\bigr]\;\mathbf 1\bigl[b\in\Av(N(t))\bigr]\,dt
  =\int_0^{\tau}\mathbf 1[\lambda_b\le t]\,dt=(\tau-\lambda_b)_+=a_b.\qedhere
\]
\end{proof}

\begin{lemma}
\label{lem:avoidable}
Let $N$ be a matroid on $E$ and $\nu=\nu(N)$.
\begin{enumerate}
  \item $|Q|-r_N(E_Q)\le k-\nu$ for every $Q\subseteq[k]$.
  \item $|\Av(N)|-r_N(E_{\Av(N)})=k-\nu$.
  \item $\Av(N)\subseteq Q$ for every $Q$ attaining equality in item 1.
\end{enumerate}
\end{lemma}

\begin{proof}
\emph{Step 1: item 1.}  A rainbow independent set $J$ with $|J|=\nu$ satisfies
\[
  \nu=|J\cap E_Q|+|J\setminus E_Q|\le r_N(E_Q)+(k-|Q|).
\]

\emph{Step 2: $|A'|=(k-\nu)+|I\cap Z|$.}  Fix a rainbow independent set $I$ with
$|I|=\nu$, and let $D_I$ be the directed graph on $E$ with arcs
\[
  y\to x\ \ \text{if }I-x+y\text{ is independent},\quad
  x\to y\ \ \text{if }\blk(y)=\blk(x)\quad(x\in I,\ y\notin I),
\]
sources the $y\notin I$ in blocks missed by $I$, and sinks the $y\notin I$ with
$I+y$ independent.  As $I$ is maximum, no path leads from a source to a
sink~\cite{lawler75,schrijver03}.  Let $Z$ be the set reachable from the sources
and $A'=\{b:b\text{ missed by }I\}\cup\blk(I\cap Z)$, so
$|A'|=(k-\nu)+|I\cap Z|$.

\emph{Step 3: $A'\subseteq\Av(N)$.}  Blocks missed by $I$ are avoidable.  For
$x\in I\cap Z$ and a shortest path $y_0\to x_1\to y_1\to\cdots\to y_{j-1}\to x_j=x$
from a source $y_0$, the arcs give $\blk(y_i)=\blk(x_i)$ for $1\le i<j$, and
$\blk(y_0)\notin\blk(I)$, so $J=I-\{x_1,\ldots,x_j\}+\{y_0,\ldots,y_{j-1}\}$ has
\[
  |J|=\nu,\qquad\blk(J)=\bigl(\blk(I)-\blk(x)\bigr)+\blk(y_0),\qquad|\blk(J)|=\nu.
\]
Hence $J$ is rainbow and misses $\blk(x)$, and it is independent by the exchange
lemma for shortest paths~\cite{schrijver03}.

\emph{Step 4: $r_N(E_{A'})\le|I\cap Z|$.}  Let $y\in E_{A'}\setminus I$.  Either
$y$ is a source or $x\to y$ for some $x\in I\cap Z$, so $y\in Z$ is not a sink,
and $I+y$ contains a circuit $C_y$.  For $x'\in C_y-y$,
\[
  I-x'+y\text{ independent}\ \Rightarrow\ y\to x'\ \Rightarrow\ x'\in I\cap Z,
\]
so $y\in\cl_N(C_y-y)\subseteq\cl_N(I\cap Z)$.  As $I$ is rainbow and misses the
blocks of $A'$ outside $\blk(I\cap Z)$, $I\cap E_{A'}=I\cap Z$, and
\[
  E_{A'}\subseteq\cl_N(I\cap Z),\qquad r_N(E_{A'})\le|I\cap Z|.
\]

\emph{Step 5: item 3.}  Let $Q$ attain equality in item 1, and let
$b\in\Av(N)\setminus Q$ be missed by a rainbow independent $J$ with $|J|=\nu$.
As $J\setminus E_Q$ is rainbow and avoids the blocks of $Q+b$,
\[
  k-|Q|=\nu-|Q|+(k-\nu)=\nu-r_N(E_Q)\le\nu-|J\cap E_Q|=|J\setminus E_Q|\le k-|Q|-1,
\]
a contradiction.

\emph{Step 6: item 2.}  By item 1 and Steps 2 and 4,
\[
  k-\nu\ge|A'|-r_N(E_{A'})\ge(k-\nu)+|I\cap Z|-|I\cap Z|=k-\nu,
\]
so $A'$ attains equality, $\Av(N)\subseteq A'$ by Step 5, and $A'=\Av(N)$ by
Step 3.
\end{proof}

\begin{proof}[Proof of Theorem~\ref{thm:validity}]
\emph{Step 1: $\varphi(t)=\Psi(p)+\int_0^{t}\varphi'(\sigma^{+})\,d\sigma$.}  For $t\in[0,\tau]$ put $a(t)_b=(t-\lambda_b)_+\in[0,t]$,
$v(t)=v-a(t)_{\blk}$ and $\varphi(t)=\Psi(p+a(t))$, so $a(0)=0$ and $a(\tau)=a$.
As a minimum of finitely many continuous piecewise linear functions of $t$,
$\varphi$ is continuous and piecewise linear, and
\[
  \varphi(t)-\varphi(0)=\int_0^{t}\varphi'(\sigma^{+})\,d\sigma.
\]

\emph{Step 2: $\varphi'(t^{+})=|S_t|-\max_{B\in\Bases(M^{v(t)})}|B\cap E_{S_t}|$.}  Let $S_t=\{b:\lambda_b\le t\}$.  For small
$h>0$, $a(t+h)=a(t)+h\mathbf 1_{S_t}$, and the minimum bases of
$v(t)-h\mathbf 1_{E_{S_t}}$ are the minimum bases of $v(t)$ with the most elements
in $E_{S_t}$.  Hence
\[
  \begin{aligned}
  \varphi'(t^{+})&=|S_t|+\frac{d}{dh}\,\beta_M\bigl(v(t)-h\mathbf 1_{E_{S_t}}\bigr)\Big|_{h=0^{+}}\\
  &=|S_t|-\max\bigl\{|B\cap E_{S_t}|:B\in\Bases(M^{v(t)})\bigr\}.
  \end{aligned}
\]

\emph{Step 3: $\varphi(t)\ge\Psi(p)+(k-r)t$.}  Let $t<\tau$.  By Lemma~\ref{lem:path} with
$y=a(t)$ and by (F3),
\[
  \Bases(M^{v(t)})\subseteq\Pos_t\subseteq\Bases(N(t)),
\]
and $\nu(N(t))=r$, $S_t=\Av(N(t))$ by Lemma~\ref{lem:updint}.  So Lemma~\ref{lem:avoidable}(2) gives, for every
$B\in\Bases(M^{v(t)})$,
\[
  |B\cap E_{S_t}|\le r_{N(t)}(E_{S_t})=|S_t|-(k-r),
\]
hence $\varphi'(t^{+})\ge k-r$ by Step 2, and by Step 1
\[
  \varphi(t)=\Psi(p)+\int_0^{t}\varphi'(\sigma^{+})\,d\sigma\ge\Psi(p)+(k-r)t\qquad(t\in[0,\tau]).
\]

\emph{Step 4: $\varphi(t)\le\Psi(p)+(k-r)t$.}  Let $B^{*}$ be a base of $M^{v}$ meeting exactly $r$
blocks.  Its $k-r$ missed blocks lie in $S_0\subseteq S_t$, as $\lambda_b=0$ on
$S_0$, and $B^{*}\in\Bases(N(t))$ by (F1) and (F2), so for $t<\tau$
\[
  |S_t|-(k-r)\le|B^{*}\cap E_{S_t}|\le r_{N(t)}(E_{S_t})=|S_t|-(k-r).
\]
Hence $v(t)(B^{*})+\sum_ba(t)_b$ has right derivative
$-(|S_t|-(k-r))+|S_t|=k-r$, so
$v(t)(B^{*})=\beta_M(v)-\sum_ba(t)_b+(k-r)t$, and
\[
  \varphi(t)\le\sum_{b=1}^{k}\bigl(p_b+a(t)_b\bigr)+v(t)(B^{*})
  =\sum_{b=1}^{k}p_b+\beta_M(v)+(k-r)t=\Psi(p)+(k-r)t.
\]

\emph{Step 5: items 1 and 2.}  By Steps 3 and 4, for $t\in[0,\tau]$,
\[
  \varphi(t)=\Psi(p)+(k-r)t,\qquad v(t)(B^{*})=\beta_M(v(t)).
\]
At $t=\tau$ these are items 1 and 2, and $\nu_0$ does not decrease since $B^{*}$
meets $r$ blocks.
\end{proof}

\subsection{Single Linkage on Possibly Optimal Elements}
\label{app:posl}

\begin{proof}[Proof of Proposition~\ref{prop:posl}]
\emph{Step 1: (F1).}  Every minimum base $B$ of $v$ lies in $\PO_0$: for $f\in B$,
Lemma~\ref{lem:minbases}(2) gives $\cl\{g:v_g<v_f\}=\cl(B\cap\{v<v_f\})\subseteq\cl(B-f)$,
which misses $f$.  At scale $0$ clusters are classes of equal values, so within a
component
\[
  u_0(x)\le u_0(y)\iff v_x\le v_y\qquad(x,y\in\PO_0).
\]
If $B$ is minimum for $v$, then $B\subseteq\PO_0$, and for $e\notin B$ and
$f\in C(B,e)$, where $v_f\le v_e$ by Lemma~\ref{lem:minbases}(1),
\[
  u_0(f)\le u_0(e)\quad\text{since}\quad
  \begin{cases}
    u_0(e)=m+1,&e\notin\PO_0,\\
    v_f\le v_e,&e\in\PO_0,
  \end{cases}
\]
so $B$ is minimum for $u_0$ by Lemma~\ref{lem:minbases}(1).  Conversely let $B$ be
minimum for $u_0$, $c\in\R$, and $K_x$ the component of $x$.  For
$x\in\PO_0\setminus B$ with $v_x\le c$, Lemma~\ref{lem:minbases}(1) gives
\[
  C(B,x)-x\subseteq\{f\in\PO_0\cap K_x:u_0(f)\le u_0(x)\}\subseteq\{v\le c\},
\]
so, by induction on $v_x$, every $x$ with $v_x\le c$ satisfies
\[
  x\in\begin{cases}
    B\cap\{v\le c\},&x\in B,\\
    \cl\bigl(C(B,x)-x\bigr)\subseteq\cl\bigl(B\cap\{v\le c\}\bigr),&x\in\PO_0\setminus B,\\
    \cl\{g:v_g<v_x\}\subseteq\cl\bigl(B\cap\{v\le c\}\bigr),&x\notin\PO_0,
  \end{cases}
\]
and $B$ is minimum for $v$ by Lemma~\ref{lem:minbases}(2).

\emph{Step 2: (F3).}  Let $\theta(B)\le t$, so that $B$ is minimum for
$u=v-t\mathbf 1_B$ by Proposition~\ref{prop:pos}.  For $f\in B$, as
$u_g\le v_g<v_f-t=u_f$ whenever $v_g<v_f-t$, Lemma~\ref{lem:minbases}(2) for $u$
gives
\[
  \cl\{g:v_g<v_f-t\}\subseteq\cl\{g:u_g<u_f\}=\cl\bigl(B\cap\{u<u_f\}\bigr)\subseteq\cl(B-f),
\]
which misses $f$, so $f\in\PO_t$.  For $e\notin B$ and $f\in C(B,e)-e$,
\[
  \begin{gathered}
  u_t(f)\le u_t(e)\quad\text{since}\\
  \begin{cases}
    u_t(e)=m+1,&e\notin\PO_t,\\
    \text{clusters are ordered by value},&e\in\PO_t,\ v_f\le v_e,\\
    0<v_f-v_e\le\theta(B)\le t\text{ links }e\text{ and }f,&e\in\PO_t,\ v_f>v_e,
  \end{cases}
  \end{gathered}
\]
and $B$ is minimum for $u_t$ by Lemma~\ref{lem:minbases}(1).

\emph{Step 3: (F2).}  A minimum base $B$ of $u_t$ lies in $\PO_t$: as
$\PO_t\supseteq\PO_0$ spans $E$, Lemma~\ref{lem:minbases}(2) with $c=m$ gives
\[
  E\subseteq\cl(\PO_t)=\cl\bigl(B\cap\{u_t\le m\}\bigr)=\cl(B\cap\PO_t),
\]
so $B\cap\PO_t=B$.  Let $t\le t'$, $B$ minimum for $u_t$, $e\notin B$ and
$f\in C(B,e)-e$, in a component $K$.  As $\PO_t\subseteq\PO_{t'}$ and linkage is
monotone, each cluster at scale $t$ lies in one at scale $t'$, in the same order.
We show $u_{t'}(f)\le u_{t'}(e)$, so that $B$ is minimum for $u_{t'}$ by
Lemma~\ref{lem:minbases}(1).  Unless $e\in\PO_{t'}\setminus\PO_t$ and $v_f>v_e$,
\[
  \begin{gathered}
  u_{t'}(f)\le u_{t'}(e)\quad\text{since}\\
  \begin{cases}
    u_{t'}(e)=m+1,&e\notin\PO_{t'},\\
    u_t(f)\le u_t(e),&e\in\PO_t,\\
    \text{clusters are ordered by value},&e\in\PO_{t'}\setminus\PO_t,\ v_f\le v_e.
  \end{cases}
  \end{gathered}
\]
Let $e\in\PO_{t'}\setminus\PO_t$ and $v_f>v_e$.  The component $K$ contains the
circuit $C(B,e)\ni f\ne e$, so it has no loop.  For $x\in\PO_t\cap K$ and $g\in K$
let
\[
  B_{\le x}=\{h\in B\cap K:u_t(h)\le u_t(x)\},\qquad H_g=\{h\in K:v_h<v_g-t\}.
\]
By induction on $v_g$, every $g\in K$ lies in $\cl(B_{\le x_g})$ for some
$x_g\in\PO_t\cap K$ with $v_{x_g}\le v_g$.  If $g\in\PO_t$, take $x_g=g$, by
Lemma~\ref{lem:minbases}(1).  If $g\notin\PO_t$, then $g\in\cl H_g$ by
Lemma~\ref{lem:minbases}(4), so $H_g\ne\varnothing$ as $g$ is not a loop; take
$x_g=x_{h_0}$ for an $h_0\in H_g$ with $u_t(x_{h_0})$ largest, so that, by the
induction hypothesis,
\[
  g\in\cl H_g\subseteq\cl(B_{\le x_g}),\qquad v_{x_g}\le v_{h_0}<v_g-t.
\]
For $x^{*}=x_e$ this gives $f\in C(B,e)-e\subseteq B_{\le x^{*}}$ and
$v_{x^{*}}<v_e-t$.  As clusters at scale $t$ lie in clusters at scale $t'$, and at
both scales clusters are ordered by value,
\[
  u_t(f)\le u_t(x^{*}),\ v_{x^{*}}<v_e<v_f\ \Rightarrow\ u_t(f)=u_t(x^{*})\
  \Rightarrow\ u_{t'}(f)=u_{t'}(x^{*})=u_{t'}(e).
\]

\emph{Step 4: (F4) and (F5).}  For $t\ge\max v-\min v$, $\{g:v_g<v_f-t\}=\varnothing$
for all $f$, so $\PO_t$ is the set of non-loops, the non-loops of each component
form one cluster, every base is minimum for $u_t$, and $N(t)=M$.  The sets
$\{g:v_g<v_f-t\}$ and the linkage are constant on each half-open interval between
consecutive levels of $\{0\}\cup\{|v_e-v_f|:e,f\in E\}$ and beyond the last one.

\emph{Step 5: $N(t)$ is determined by $\mathcal G_t$.}  We have
\[
  v_g<v_f-t\iff(g,f)\notin\mathcal G_t,\qquad |v_x-v_y|\le t\iff(x,y),(y,x)\in\mathcal G_t,
\]
and for $x,y$ in different clusters of a component $v_x<v_y\iff(x,y)\notin\mathcal G_t$;
the components depend on $M$ alone.
\end{proof}

\subsection{Breakpoints}
\label{app:events}

\begin{proof}[Proof of Lemma~\ref{lem:events}]
\emph{Step 1: item 1.}  A loop lies in $\cl\varnothing$, hence in no $\PO_t$.
Let $x$ be a non-loop and $0=d_0\le\cdots\le d_m$ the entries of $D^{x}$ in
increasing order with a zero prepended.  For $t\ge0$, $v_g<v_x-t$ if and only if
$D^{x}_g>t$, so
\[
  x\notin\PO_t\iff x\in\cl\bigl(E\setminus G_t(D^{x})\bigr)\iff\phi_x\bigl(G_t(D^{x})\bigr)=1.
\]
The function $\phi_x$ is antitone, as $\cl$ is monotone, and $\phi_x(E)=0$, as $x$ is
not a loop.  Let $j$ be the least index with $\phi_x(G_{d_j}(D^{x}))=0$; it exists,
as $G_{d_m}(D^{x})=E$.  The sets $G_{d_i}(D^{x})$ grow with $i$, so
$\phi_x(G_{d_i}(D^{x}))=\mathbf 1[i<j]$, and $d_i<d_j$ for $i<j$, since equal values
give equal sets.  Hence
\[
  t_x=\sum_{i=0}^{m-1}\Delta_i\,\mathbf 1[i<j]=d_j,
\]
and as $G_t(D^{x})=G_{d_i}(D^{x})$ for the largest $i$ with $d_i\le t$,
\[
  x\notin\PO_t\iff i<j\iff t<d_j=t_x.
\]

\emph{Step 2: item 2.}  For joined $a,b$ and $t\ge0$, the link of
Definition~\ref{def:posl} at scale $t$ asks $a,b\in\PO_t$ and $|v_a-v_b|\le t$, and
for $a=b$ only $a\in\PO_t$.  By Step 1,
\[
  a,b\text{ linked at scale }t\iff W_{ab}\le t,
\]
and an element of $\PO_t$ is a non-loop, linked only to joined elements.  As the
clusters at scale $t$ are the classes of the chains of links,
\[
  \begin{aligned}
  x,y\text{ in one cluster}
  &\iff\max_iW_{a_ia_{i+1}}\le t\text{ for some sequence }(a_i)\text{ of item 2}\\
  &\iff T_{xy}\le t.
  \end{aligned}
\]

\emph{Step 3: item 3.}  Concatenating sequences gives
$T_{xz}\le\max(T_{xy},T_{yz})$, reversing them gives $T_{xy}=T_{yx}$, and as
$W_{ab}\ge\max(t_a,t_b)$ and the sequence $x,x$ is allowed,
\[
  \max(t_x,t_y)\le T_{xy},\qquad T_{xx}=t_x.
\]
Let $h=T_{xy}$ and $C=\{z:z\text{ joined to }x,\ T_{xz}\le h\}$, which contains $x$
and $y$.  If $t_z=h$ for some $z\in C$, then $h\in\Lambda$.  Otherwise $t_z<h$ for
every $z\in C$, as $t_z\le T_{xz}\le h$.  Let $z_0$ be the least element of $C$.  As
$h=T_{xy}\le\max(T_{xz_0},T_{z_0y})$, some $e\in\{x,y\}$ has $T_{z_0e}\ge h$.  The set
$P=\{z\in C:T_{ez}<h\}$ contains $e$, as $T_{ee}=t_e<h$, and misses $z_0$; let $z$ be
its least element, so $z_0<z$.  For $y'<z$ joined to $z$ with $T_{zy'}<h$,
\[
  T_{xy'}\le\max(T_{xz},T_{zy'})\le h,\qquad T_{ey'}\le\max(T_{ez},T_{zy'})<h,
\]
so $y'\in P$, which contradicts the choice of $z$.  So $T_{zy'}\ge h$ for every
$y'<z$ joined to $z$, and $y'=z_0$ attains $h$:
\[
  h\le T_{zz_0}\le\max(T_{zx},T_{xz_0})\le h.
\]
Hence $\sigma_z=h\in\Lambda$.

\emph{Step 4: item 4.}  Let $x,y$ be joined and $t\ge0$.  By
Definition~\ref{def:posl}, $u_t(x)<u_t(y)$ if and only if $x\in\PO_t$ and either
$y\notin\PO_t$, or $x,y$ lie in different clusters with $v_x<v_y$, which then means
$v_y-v_x>t$.  By Steps 1 and 2,
\[
  \begin{aligned}
  u_t(x)<u_t(y)&\iff t\ge t_x\ \text{and}\
  \bigl(t<t_y\ \text{or}\ t<\min(T_{xy},v_y-v_x)\bigr)\\
  &\iff t_x\le t<\omega_{xy},
  \end{aligned}
\]
as $t\ge t_x$ makes $t<\max(t_x,c)$ equivalent to $t<c$.  By the definition of
$\Lambda$ and Step 3,
\[
  \omega_{xy}=\begin{cases}
    \max(t_x,t_y)\in\Lambda,&\min(T_{xy},v_y-v_x)\le\max(t_x,t_y),\\
    T_{xy}\in\Lambda\cup\{0\},&\max(t_x,t_y)<T_{xy}\le v_y-v_x,
  \end{cases}
\]
and the remaining case $\max(t_x,t_y)<v_y-v_x<T_{xy}$ does not occur: it forces
$x\ne y$, and then
\[
  T_{xy}\le W_{xy}=\max\bigl(t_x,t_y,|v_x-v_y|\bigr)=v_y-v_x.
\]

\emph{Step 5: $N(t)$ is constant between consecutive values of $\Lambda\cup\{0\}$.}  By Steps 1, 2 and 4, the bits $\mathbf 1[t_x\le t]$,
$\mathbf 1[T_{xy}\le t]$ and $\mathbf 1[\omega_{xy}\le t]$ determine $\PO_t$, the
clusters and their order, hence $u_t$ and $N(t)$, and by Steps 3 and 4 each of these
bits changes only at a value of $\Lambda\cup\{0\}$.
\end{proof}

\subsection{Termination}
\label{app:term}

\begin{proof}[Proof of Theorem~\ref{thm:potential}]
\emph{Step 1: item 1.}  For $b\in S_0$, some base of $M^{v}$ meeting exactly $r$
blocks misses $b$; by Theorem~\ref{thm:validity}(2) it is a base of $M^{v'}$, so
$b\in\Av(M^{v'})$ when $\nu(M^{v'})=r$.

\emph{Step 2: item 2.}  We have $a_b=\tau$ for $b\in S_0$, where $\lambda_b=0$,
and $a_b=(\tau-\lambda_b)_+\le\tau$ otherwise, so for $y\in E_{S_0}$ and
$g\notin E_{S_0}$
\[
  v'_y-v'_g=(v_y-\tau)-\bigl(v_g-a_{\blk(g)}\bigr)\le v_y-v_g.
\]

\emph{Step 3: $0<\delta\le\tau$, and $a_b\le\tau-\delta$ for $b\notin S_0$.}  Let
$\delta=\min\{t:\Av(N(t))\ne S_0\text{ or }\nu(N(t))>r\}$, which exists by (F4)
and (F5).  Since $N(0)=M^{v}$ and $\nu(N(\tau))\ge r+1$, $0<\delta\le\tau$, and
since $\Av(N(t))=S_0$ for $t<\delta$, every $b\notin S_0$ has
\[
  \lambda_b\ge\delta,\qquad a_b\le\tau-\delta.
\]

\emph{Step 4: $r_{N(\delta)}(E_{S_0})>|S_0|-(k-r)$.}  By Lemma~\ref{lem:avoidable}(2) for $M^{v}$,
$r_{M^{v}}(E_{S_0})=|S_0|-(k-r)$.  If $\nu(N(\delta))=r$, then
$S_0\subsetneq\Av(N(\delta))$ by (F1), (F2) and the choice of $\delta$, so by
Lemma~\ref{lem:avoidable}(3) the set $S_0$ does not attain equality in
Lemma~\ref{lem:avoidable}(1).  As $\nu(N(\delta))\ge r$, Lemma~\ref{lem:avoidable}(1)
gives
\[
  |S_0|-r_{N(\delta)}(E_{S_0})
  \begin{cases}
    \le k-\nu(N(\delta))<k-r,&\nu(N(\delta))>r,\\
    <k-\nu(N(\delta))=k-r,&\nu(N(\delta))=r,
  \end{cases}
\]
so
\begin{equation}
  r_{N(\delta)}(E_{S_0})>|S_0|-(k-r).
  \label{eq:jump}
\end{equation}

\emph{Step 5: $g\in B\setminus E_{S_0}$ and $y\in E_{S_0}\setminus B$ in one cluster at scale $\delta$ with $v_g<v_y$.}  Let $B$ be a base of $M^{v}$ with
$|B\cap E_{S_0}|=r_{M^{v}}(E_{S_0})$.  By (F1) and (F2), $B\in\Bases(N(\delta))$, and
by \eqref{eq:jump},
\[
  |B\cap E_{S_0}|=|S_0|-(k-r)<r_{N(\delta)}(E_{S_0}),
\]
so $B\cap E_{S_0}$ does not span $E_{S_0}$ in $N(\delta)$, and some
$y\in E_{S_0}\setminus B$ has, in $N(\delta)$, a circuit in $B+y$ through some
$g\in B\setminus E_{S_0}$.  By Lemma~\ref{lem:minbases}(3) for $u_\delta$ and
$B\subseteq\PO_\delta$ (Step 3 of the proof of Proposition~\ref{prop:posl}),
\[
  g\in C(B,y),\qquad u_\delta(g)=u_\delta(y)\le m,
\]
so $g$ and $y$ lie in one cluster at scale $\delta$.  By Lemma~\ref{lem:minbases}(1)
for $v$, $v_g\le v_y$.  If $v_g=v_y$, then $B-g+y$ is a base of $M^{v}$ by
Lemma~\ref{lem:minbases}(3), and
\[
  r_{M^{v}}(E_{S_0})+1=|(B-g+y)\cap E_{S_0}|\le r_{M^{v}}(E_{S_0}),
\]
a contradiction.  So $v_g<v_y$.

\emph{Step 6: $0<\beta-\alpha\le\delta$.}  In this cluster let
\[
  \alpha=\max\{v_z:z\notin E_{S_0},\ v_z<v_y\}=v_{z_1},\quad
  \beta=\min\{v_z:z\in E_{S_0},\ v_z>\alpha\}=v_{z_2},
\]
the sets being nonempty by $g$ and $y$, so $\beta\le v_y$.  By the choice of
$\alpha$ and $\beta$ no value of the cluster lies in $(\alpha,\beta)$, so they are
consecutive and $0<\beta-\alpha\le\delta$.

\emph{Step 7: item 3.}  By Steps 2 and 3,
\[
  v'_{z_2}-v'_{z_1}\le(\beta-\tau)-\bigl(\alpha-(\tau-\delta)\bigr)=(\beta-\alpha)-\delta\le0,
\]
so $(z_2,z_1)\in\Pi(v')\setminus\Pi(v)$, as $\beta>\alpha$.

\emph{Step 8: at most $r(\lfloor m^{2}/4\rfloor+1)$ updates.}  By Lemma~\ref{lem:avoidable}(2) for $M^{v}$ and
$r_{M^{v}}(E)=k$,
\[
  |S_0|-r_{M^{v}}(E_{S_0})=k-r\ge1,\qquad\text{so}\qquad k-r\le|S_0|\le k-1.
\]
By item 1, $S_0$ only grows within a stage, so it takes at most $r$ values, and
while it is fixed, items 2 and 3 make $\Pi$ grow strictly at every update, with
\[
  |\Pi|\le|E_{S_0}|\,|E\setminus E_{S_0}|\le\lfloor m^{2}/4\rfloor.
\]
So a stage has at most $r(\lfloor m^{2}/4\rfloor+1)$ updates.
\end{proof}

\begin{lemma}
\label{lem:crossing}
Let $B$ be a base of $M^{v}$ meeting exactly $r$ blocks, and for $u\in\R^{E}$ let
\[
  R_B(u)=\bigl\{(x,z)\in E\times E:x\ne z,\ x\in B\text{ or }z\in B,\
  \blk(x)\in\Av(M^{u}),\ u_x\le u_z\bigr\}.
\]
\begin{enumerate}
  \item Some $x\in E_{S_0}$ and $z\notin E_{S_0}$ in the same connected component of $M$,
        one of them in $B$, satisfy $v_x>v_z$ and $v'_x\le v'_z$.
  \item If $|S_0|=k-r$ and $\nu(M^{v'})=r$, then $\Av(M^{v'})\ne S_0$.
  \item If $\nu(M^{v'})=r$, then $R_B(v)\subsetneq R_B(v')$.
\end{enumerate}
\end{lemma}

\begin{proof}
We use $\delta$ of Step 3 of the proof of Theorem~\ref{thm:potential} and write
$\rho=|S_0|-(k-r)$.

\emph{Step 1: $|B\cap E_{S_0}|=\rho$.}  The $k-r$ blocks that $B$ misses lie in $S_0$, so
$B\cap E_{S_0}$ meets at least $\rho$ blocks, and it is independent in $M^{v}$; by
Lemma~\ref{lem:avoidable}(2),
\[
  \rho\le|B\cap E_{S_0}|\le r_{M^{v}}(E_{S_0})=\rho.
\]
So Step 5 of the proof of Theorem~\ref{thm:potential} applies to $B$: it gives
$g\in B\setminus E_{S_0}$ and $y\in E_{S_0}\setminus B$ in one cluster $\Gamma$ at
scale $\delta$ with $v_g<v_y$, and its last argument gives, for
$f\in E_{S_0}\setminus B$,
\[
  v_c<v_f\qquad\bigl(c\in C(B,f)\setminus E_{S_0}\bigr).
\]

\emph{Step 2: $0<v_x-v_z\le\delta$ for some $x\in E_{S_0}$ and $z\notin E_{S_0}$, one of them in $B$.}  Let
\[
  v_{z_2}=\min\{v_z:z\in\Gamma\cap E_{S_0},\ v_z>v_g\},\qquad
  v_{z_1}=\max\{v_z:z\in\Gamma\setminus E_{S_0},\ v_z<v_{z_2}\},
\]
the sets being nonempty by $y$ and $g$.  Then $v_g\le v_{z_1}<v_{z_2}$, no value
of $\Gamma$ lies in $(v_{z_1},v_{z_2})$, and as in Step 6 of the proof of
Theorem~\ref{thm:potential}, $v_{z_2}-v_{z_1}\le\delta$.  We find $x\in E_{S_0}$ and
$z\notin E_{S_0}$, one of them in $B$, with
\begin{equation}
  0<v_x-v_z\le\delta.
  \label{eq:pair}
\end{equation}
If $z_2\in B$, take $(x,z)=(z_2,z_1)$.  Otherwise $z_2\in\PO_\delta$ lies in
$\cl(C(B,z_2)-z_2)$ but not in $\cl\{f:v_f<v_{z_2}-\delta\}$, so by the minimality
of $B$ some $c\in C(B,z_2)-z_2$ has
\[
  v_{z_2}-\delta\le v_c\le v_{z_2}.
\]
As $c\in B\subseteq\PO_\delta$ lies in the component of $z_2$ and within $\delta$
of it, $c\in\Gamma$, and \eqref{eq:pair} holds, with $x$ and $z$ in the cluster
$\Gamma$ and so in one component, for
\[
  (x,z)=
  \begin{cases}
    (z_2,c),&c\notin E_{S_0},\ \text{where }v_c<v_{z_2}\text{ by Step 1},\\
    (c,z_1),&c\in E_{S_0},\ v_c=v_{z_2},\\
    (z_2,g),&c\in E_{S_0},\ v_c<v_{z_2},
  \end{cases}
\]
in the last case because $v_c\le v_g$ by the choice of $z_2$, so that
$0<v_{z_2}-v_g\le v_{z_2}-v_c\le\delta$.

\emph{Step 3: item 1.}  By Steps 2 and 3 of the proof of
Theorem~\ref{thm:potential} and \eqref{eq:pair},
\[
  v'_x-v'_z\le(v_x-\tau)-\bigl(v_z-(\tau-\delta)\bigr)=(v_x-v_z)-\delta\le0.
\]

\emph{Step 4: item 2.}  Let $\rho=0$, and suppose $\nu(M^{v'})=r$ and
$\Av(M^{v'})=S_0$.  By \eqref{eq:jump}, $r_{N(\delta)}(E_{S_0})>0$, so some
$y\in E_{S_0}$ lies in a base of $N(\delta)$, hence in $\PO_\delta$.  By
Lemma~\ref{lem:avoidable}(2) for $M^{v'}$, $r_{M^{v'}}(E_{S_0})=\rho=0$, so a
minimum base $B'$ of $v'$ misses $E_{S_0}$, and the last argument of Step 1, at
$v'$ and $B'$, gives $v'_f<v'_y$ for $f\in C(B',y)-y$.  By Step 3 of the proof of
Theorem~\ref{thm:potential},
\[
  v_f\le v'_f+(\tau-\delta)<v'_y+(\tau-\delta)=v_y-\delta\qquad\bigl(f\in C(B',y)-y\bigr),
\]
so $y\in\cl\bigl(C(B',y)-y\bigr)\subseteq\cl\{f:v_f<v_y-\delta\}$, against
$y\in\PO_\delta$.

\emph{Step 5: item 3.}  Let $\nu(M^{v'})=r$ and $(x,z)\in R_B(v)$.  By
Theorem~\ref{thm:potential}(1), $\blk(x)\in S_0\subseteq\Av(M^{v'})$, and by Step 2
of its proof, $a_{\blk(x)}=\tau\ge a_{\blk(z)}$, so
\[
  v'_x-v'_z=(v_x-\tau)-\bigl(v_z-a_{\blk(z)}\bigr)\le v_x-v_z\le0,
\]
and $(x,z)\in R_B(v')$.  The pair $(x,z)$ of item 1 has $x\ne z$, as $\blk(x)\in S_0$
and $\blk(z)\notin S_0$; it lies outside $R_B(v)$, as $v_x>v_z$, and in $R_B(v')$,
as $\blk(x)\in S_0\subseteq\Av(M^{v'})$ and $v'_x\le v'_z$.
\end{proof}

\begin{proof}[Proof of Theorem~\ref{thm:stagelength}]
Write $\rho=|S_0|-(k-r)$.

\emph{Step 1: at most $1+(r-1)km$ updates.}  Call a maximal run of updates of stage $r$
that start
from states with the same $S_0$ a \emph{phase}.  By Theorem~\ref{thm:potential}(1)
and Step 8 of its proof, $S_0$ grows strictly from one phase to the next and
$0\le\rho\le r-1$, so the phases have distinct values of $\rho$, and by Lemma~\ref{lem:crossing}(2) the
phase with $\rho=0$, if any, has one update.  In a phase with $\rho\ge1$ fix a
base $B$ of $M^{v}$ meeting exactly $r$ blocks at its first state; by Theorem~\ref{thm:validity}(2) it is a base of $M^{v}$
meeting $r$ blocks at every state of the phase.  Let
\[
  \Phi(v)=\bigl|\{(x,z):x\in E_{S_0},\ z\notin E_{S_0},\ x\in B\text{ or }z\in B,\
  v_x>v_z\}\bigr|,
\]
formed with the $S_0$ of the phase.  At an update $v\to v'$ of the phase, by Step 2
of the proof of Theorem~\ref{thm:potential} no pair with $v_x\le v_z$ has
$v'_x>v'_z$, and by Lemma~\ref{lem:crossing}(1) some pair counted by $\Phi(v)$ has $v'_x\le v'_z$, so
$\Phi(v')\le\Phi(v)-1$, while
\[
  \Phi(v)\le|B\cap E_{S_0}|\,|E\setminus E_{S_0}|+|E_{S_0}|\,|B\setminus E_{S_0}|\le|B|\,m=km.
\]
So a phase with $\rho\ge1$ has at most $km$ updates, and stage $r$ has at most
$1+(r-1)km$.

\emph{Step 2: at most $1+k(2m-k-1)$ updates.}  Let stage $r$ consist of the updates
$v^{j}\to v^{j+1}$, $j<n$, so that $\nu(M^{v^{j}})=r$ for $j<n$, and let $B$ be a base
of $M^{v^{0}}$ meeting exactly $r$ blocks.  By Theorem~\ref{thm:validity}(2) and
induction on $j$, $B$ is a base of $M^{v^{j}}$ meeting exactly $r$ blocks for every
$j<n$, so Lemma~\ref{lem:crossing}(3) applies with this $B$ at every update $v^{j}\to v^{j+1}$ with
$j<n-1$.  As $R_B(v^{n-1})$ consists of pairs $(x,z)$ with $x\ne z$ and $x$ or $z$
in $B$, and $|B|=k$,
\[
  n-1\le\bigl|R_B(v^{n-1})\bigr|-\bigl|R_B(v^{0})\bigr|\le k(m-1)+(m-k)k=k(2m-k-1).
\]
With Step 1, stage $r$ has at most $\min\bigl(1+(r-1)km,\,1+k(2m-k-1)\bigr)$ updates.
\end{proof}

\begin{proof}[Proof of Corollary~\ref{cor:rankone}]
\emph{Step 1: $\{x,z\}$ is a circuit.}  Let $B$, $x$ and $z$ be as in
Lemma~\ref{lem:crossing}(1).  Since $x\in E_{S_0}$ and $z\notin E_{S_0}$, $x\ne z$, and
as their component has rank one, $\{x,z\}$ is a circuit.

\emph{Step 2: $z\in B$ and $x\notin B$.}  If $x\in B$, then $z\notin B$ and $C(B,z)=\{x,z\}$, and
Lemma~\ref{lem:minbases}(1) would give
\[
  v_x\le v_z<v_x.
\]
So $z\in B$, $x\notin B$ and $C(B,x)=\{x,z\}$.

\emph{Step 3: $\nu(M^{v'})>r$ or $\Av(M^{v'})\supsetneq S_0$.}  By Theorem~\ref{thm:validity}(2), $B$ is a minimum base
of $v'$, so by Lemma~\ref{lem:minbases}(1) and Lemma~\ref{lem:crossing}(1),
\[
  v'_z\le v'_x\le v'_z,
\]
and by Lemma~\ref{lem:minbases}(3) $B-z+x$ is a minimum base of $v'$ with
\[
  |(B-z+x)\cap E_{S_0}|=|B\cap E_{S_0}|+1=\rho+1,
\]
by Step 1 of the proof of Lemma~\ref{lem:crossing}.  If $\nu(M^{v'})=r$ and
$\Av(M^{v'})=S_0$, then Lemma~\ref{lem:avoidable}(2) for $M^{v'}$ gives
\[
  \rho+1=|(B-z+x)\cap E_{S_0}|\le r_{M^{v'}}(E_{S_0})=|S_0|-(k-r)=\rho,
\]
a contradiction.  So $\nu(M^{v'})>r$ or $\Av(M^{v'})\supsetneq S_0$ by
Theorem~\ref{thm:potential}(1).

\emph{Step 4: at most $r$ updates.}  By Step 8 of the proof of Theorem~\ref{thm:potential},
$k-r\le|S_0|\le k-1$ within the stage, and by Step 3 every update of the stage
that does not end it increases $|S_0|$.  So the stage has at most
$(k-1)-(k-r)+1=r$ updates.
\end{proof}

\subsection{The Circuit}
\label{app:circuit}

\begin{proof}[Proof of Lemma~\ref{lem:augment}]
\emph{Step 1: a round of $m^{2}q+8m^{3}+12m^{2}+12m+4$ gates enlarges $I$ unless $|I|=\nu(N_x)$.}  From the indicator vector of a rainbow independent set
$I$, a round builds the exchange graph $D_I$ of the proof of
Lemma~\ref{lem:avoidable}: the circuit of $y$ tests $I-x'+y$ for the arcs
$y\to x'$ and $I+y$ for the sinks, and the sources and the arcs inside blocks are
read from $I$.  An answer is used only when $y\notin I$, and for an arc only when
$x'\in I$; the bits of $I$ mask the others.  The used sets contain $y$, and
removing $y$ leaves $I-x'$ or $I$, which are independent, so the circuits answer
them.  With a breadth-first search from the sources, the sink of least index in
the first layer containing a sink, the predecessors of least index and the
exchange along this shortest path, a round costs
\[
  \underbrace{m(m-1)q}_{\text{arcs}}+\underbrace{mq}_{\text{sinks}}
  +\underbrace{8m^{3}+12m^{2}+12m+4}_{\text{search and exchange}}
  =m^{2}q+8m^{3}+12m^{2}+12m+4.
\]
By the exchange lemma for shortest paths~\cite{schrijver03} the new set is
rainbow, independent and has $|I|+1$ elements, and if no sink is reachable, $I$
is maximum and is kept~\cite{lawler75,schrijver03}.

\emph{Step 2: $\min(j,\nu(N_x))\le|I|\le\nu(N_x)$ after $j$ rounds.}  From $I=\varnothing$, $|I|$ grows by one per round
until $|I|=\nu(N_x)\le k$, so $j$ rounds give
\[
  \min\bigl(j,\nu(N_x)\bigr)\le|I|\le\nu(N_x),
\]
and $k$ rounds give $|I|=\nu(N_x)$.

\emph{Step 3: the bits $\mathbf 1[\nu(N_x)\le r_0]\wedge\mathbf 1[b\in\Av(N_x)]$.}  Run $j=r_0+1$ rounds.  If $\nu(N_x)\le r_0$, then
$|I|=\nu(N_x)$ by Step 2, and one more exchange graph ($m(m-1)$ tests) and its
reachable set $Z$ ($2m^{3}$ gates) give, by Steps 2 to 6 of the proof of
Lemma~\ref{lem:avoidable},
\[
  \Av(N_x)=\{b:b\text{ missed by }I\}\cup\blk(I\cap Z).
\]
If $\nu(N_x)>r_0$, then $|I|\ge r_0+1$ by Step 2.  In both cases
\[
  \begin{aligned}
  &\mathbf 1\bigl[|I|\le r_0\bigr]\wedge\mathbf 1\bigl[b\in\{b:b\text{ missed by }I\}\cup\blk(I\cap Z)\bigr]\\
  &\qquad=\mathbf 1\bigl[\nu(N_x)\le r_0\bigr]\wedge\mathbf 1\bigl[b\in\Av(N_x)\bigr].
  \end{aligned}
\]
A population count of $I$, a comparison with $r_0$ and $k$ conjunctions complete
the outputs, within
$(r_0+2)(m+1)^{2}(q+10(m+1))+3k(m+1)+5\lceil\log_2(r_0+1)\rceil$ gates.
\end{proof}

\begin{proof}[Proof of Lemma~\ref{lem:observer}]
Fix $v$ and $t\ge0$, and write $\PO=\PO_t$, $u=u_t$ and $N_t=N(t)$.

\emph{Step 1: $\PO$, the clusters and their order with $2m^{2}$ gates.}  The loops and the connected components of $M$ are
fixed, and by Lemma~\ref{lem:events}, at the bits $\mathbf 1[\mathcal Q(v)\le t]$,
\[
  x\in\PO\iff x\text{ is a non-loop and }t_x\le t,
\]
two distinct elements lie in one cluster if and only if they are joined and
$T_{xy}\le t$, and for $x,y$ in one component
\[
  u(x)<u(y)\iff x,y\text{ joined},\ t_x\le t\ \text{and not}\ \omega_{xy}\le t.
\]
So the bits of $\PO$ and of the clusters cost no gate, and the order bits cost
$2m^{2}$ gates.

\emph{Step 2: the independence test of $N_t$ with $q=s+O(m)$ gates.}  For $y\in\PO$ let $\mathrm{Lev}(y)$ be the cluster
of $y$ and $\mathrm{Below}(y)$ the set of elements of its component with smaller
$u$; both depend only on the cluster.  By \eqref{eq:rank-mu} within components,
\[
  r_{N_t}(A)=\sum_{x}\kappa_x(A),\qquad
  \kappa_x(A)=r_M\bigl(\mathrm{Below}(x)\cup(A\cap\mathrm{Lev}(x))\bigr)-r_M\bigl(\mathrm{Below}(x)\bigr),
\]
the sum over one element $x$ of each cluster; the set $E\setminus\PO=\{u=m+1\}$
contributes $0$, as $\PO\supseteq\PO_0$ spans $E$.  Each term satisfies
$\kappa_x(A)\le|A\cap\mathrm{Lev}(x)|$, and the clusters partition $\PO$, so
\[
  r_{N_t}(A)\le\sum_x|A\cap\mathrm{Lev}(x)|=|A\cap\PO|\le|A|,
\]
and $A$ is independent in $N_t$ if and only if $A\subseteq\PO$ and
$\kappa_x(A)=|A\cap\mathrm{Lev}(x)|$ for every $x$.  Let $y\in A$ with $A-y$
independent, so that $A-y\subseteq\PO$ and, for $x$ outside the cluster of $y$,
\[
  \kappa_x(A)=\kappa_x(A-y)=|(A-y)\cap\mathrm{Lev}(x)|=|A\cap\mathrm{Lev}(x)|.
\]
Hence
\[
  \begin{aligned}
  &A\text{ independent in }N_t\iff y\in\PO\ \text{and}\\
  &\qquad r_M\bigl(\mathrm{Below}(y)\cup(A\cap\mathrm{Lev}(y))\bigr)
  =r_M\bigl(\mathrm{Below}(y)\bigr)+|A\cap\mathrm{Lev}(y)|.
  \end{aligned}
\]
The ranks $r_M(\mathrm{Below}(y))$ cost $ms$ gates once, and then this test costs
$q=s+O(m)$ gates: the sets $A\cap\mathrm{Lev}(y)$ and
$\mathrm{Below}(y)\cup(A\cap\mathrm{Lev}(y))$, one rank evaluation, a subtraction,
a population count, a comparison and a conjunction with $\mathbf 1[y\in\PO]$.

\emph{Step 3: the bits $F_b$ with $O(r_0m^{2}(s+m))$ gates.}  Lemma~\ref{lem:augment}, with the input bits and
the data of Steps 1 and 2 as its input, the matroids $N_t$ and the tests of Step 2,
gives the bits $F_b$ from $r_0+1$ rounds with
\[
  (r_0+2)(m+1)^{2}\bigl(q+10(m+1)\bigr)+3k(m+1)+5\lceil\log_2(r_0+1)\rceil
  =O\bigl(r_0m^{2}(s+m)\bigr)
\]
gates, as $k\le m$, and Steps 1 and 2 cost $ms+2m^{2}$ gates besides the tests.
\end{proof}

\begin{proof}[Proof of Lemma~\ref{lem:oneupdate}]
\emph{Step 1: $v$ with $m$ gates.}  One subtraction per element gives $v=w-p_{\blk}$.

\emph{Step 2: $\Lambda$ and $\mathcal Q(v)$ with $O(m^{2}s+m^{3})$ gates.}  For a non-loop $x$, the vector $D^{x}$ costs $2m$
gates and the Boolean circuit of $\phi_x(S)=\mathbf 1[r_M(E\setminus S+x)=r_M(E\setminus S)]$,
with $m$ negations, two rank evaluations and a comparison, has $2s+O(m)$ gates, so
Theorem~\ref{thm:compiler} gives $t_x$ with $O(ms+m^{2})$ gates.  Let $\mathcal J$ be the
set of joined pairs, so that $W\in\R^{\mathcal J}$.  The Floyd--Warshall algorithm
$\mathrm{FW}$ over $(\min,\max)$ starts from $W$ and replaces $T_{xy}$ by
$\min(T_{xy},\max(T_{xz},T_{zy}))$ for the pivots $z$ in turn, and Warshall's
algorithm~\cite{warshall62} does the same on bits with $\vee$ and $\wedge$.  As
\[
  \mathbf 1\bigl[\min(T_{xy},\max(T_{xz},T_{zy}))\le t\bigr]
  =\mathbf 1[T_{xy}\le t]\vee\bigl(\mathbf 1[T_{xz}\le t]\wedge\mathbf 1[T_{zy}\le t]\bigr),
\]
every pivot commutes with the threshold map, so for every $t\ge0$
\[
  \begin{tikzcd}[column sep=4.5em,row sep=2.2em]
    \R^{\mathcal J}\arrow[r,"\mathrm{FW}"]\arrow[d,"{\mathbf 1[\,\cdot\,\le t]}"'] &
    \R^{\mathcal J}\arrow[d,"{\mathbf 1[\,\cdot\,\le t]}"]\\
    \{0,1\}^{\mathcal J}\arrow[r,"\text{Warshall}"'] & \{0,1\}^{\mathcal J}
  \end{tikzcd}
\]
commutes, and $\mathrm{FW}(W)$ is the matrix of the least values $T_{xy}$ of
Lemma~\ref{lem:events}(2).  With the
values $\sigma_x$ and $\omega_{xy}$, the levels $\Lambda$ and the values
$\mathcal Q(v)$ cost
\[
  \underbrace{O(m^{2}s+m^{3})}_{t_x}+\underbrace{5m^{2}}_{W}+\underbrace{4m^{3}}_{T}
  +\underbrace{3m^{2}}_{\sigma}+\underbrace{6m^{2}}_{\omega}=O(m^{2}s+m^{3})
\]
gates.

\emph{Step 3: $a_b=\sum_{i<2m}\Delta_i\,F_b(\mathbf 1[\mathcal Q(v)\le d_i])$.}  Let $F_b$ be the bits of
Lemma~\ref{lem:observer} with $r_0=r$, and for a matroid $N$ on $E$ let
$\chi_b(N)=\mathbf 1[\nu(N)\le r]\wedge\mathbf 1[b\in\Av(N)]$.  By
Lemma~\ref{lem:observer}, for every $v\in\R^{E}$
\[
  \begin{tikzcd}[column sep=5.4em,row sep=2.4em]
    {[0,\infty)}\arrow[r,"{t\,\mapsto\,\mathbf 1[\mathcal Q(v)\le t]}"]
      \arrow[d,"{t\,\mapsto\,N(t)}"'] &
    \{0,1\}^{m+2m^{2}}\arrow[d,"F_b"]\\
    \{\text{matroids on }E\}\arrow[r,"\chi_b"'] & \{0,1\}
  \end{tikzcd}
\]
commutes.  Let $0=d_0\le\cdots\le d_{2m}$ be the entries of
$\Lambda$ in increasing order with a zero prepended.  Every entry of $\mathcal Q(v)$
lies in $\Lambda\cup\{0\}$ by Lemma~\ref{lem:events}, so the bits
$\mathbf 1[\mathcal Q(v)\le t]$ are constant on each $[d_i,d_{i+1})$, and for
$t\ge d_{2m}$ they are all $1$, as at $t'=\max(d_{2m},\max v-\min v)$, where
$N(t')=M$ by Step 4 of the proof of Proposition~\ref{prop:posl} and $\nu(M)=k>r$.
So for $t\ge d_{2m}$, by the square,
\[
  F_b\bigl(\mathbf 1[\mathcal Q(v)\le t]\bigr)=F_b\bigl(\mathbf 1[\mathcal Q(v)\le t']\bigr)
  =\chi_b(M)=0.
\]
If $\nu_0=r$, then by Lemma~\ref{lem:updint} and the square,
\[
  \begin{aligned}
  a_b&=\int_0^{\infty}\chi_b\bigl(N(t)\bigr)\,dt
  =\int_0^{\infty}F_b\bigl(\mathbf 1[\mathcal Q(v)\le t]\bigr)\,dt\\
  &=\sum_{i=0}^{2m-1}\Delta_i\,F_b\bigl(\mathbf 1[\mathcal Q(v)\le d_i]\bigr),
  \end{aligned}
\]
and if $\nu_0>r$, then $\nu(N(t))\ge\nu_0>r$ for all $t$ and both sides are
$0=a_b$.

\emph{Step 4: $a$ with $O(rm^{3}(s+m))$ gates.}  Corollary~\ref{cor:levels} with the $2m$
levels $\Lambda$, the $m+2m^{2}$ values $\mathcal Q(v)$ as $Y$ and the $k$ outputs
$F_b$, applied to the circuit of Lemma~\ref{lem:observer}, gives $a$ with
\[
  6m^{2}+2m\Bigl(1+4(m+2m^{2})+3\cdot O\bigl(rm^{2}(s+m)\bigr)+k\Bigr)
  =O\bigl(rm^{3}(s+m)\bigr)
\]
gates.

\emph{Step 5: $p+a$ with $k$ gates.}  One addition per block gives $p+a$.
\end{proof}

\begin{proof}[Proof of Theorem~\ref{thm:C}]
\emph{Step 1: $\ell_r$ updates of stage $r$ for $r=1,\ldots,k-1$.}  Start from $p=0$, where $\nu_0\ge1$ since a minimum
base of $w$ meets some block.  For $r=1,\ldots,k-1$ apply the circuit of
Lemma~\ref{lem:oneupdate} for stage $r$, which is the identity at $\nu_0>r$ as
$a=0$ there, exactly $\ell_r$ times, where
\[
  \ell_r=\min\bigl(r(\lfloor m^{2}/4\rfloor+1),\,1+(r-1)km,\,1+k(2m-k-1)\bigr).
\]

\emph{Step 2: $\nu_0=k$ after stage $k-1$.}  By induction on $r$, $\nu_0\ge r$ at the start of
stage $r$.  As $\nu_0$ never decreases (Theorem~\ref{thm:validity}(2)) and
Theorems~\ref{thm:potential} and~\ref{thm:stagelength} bound the number of
updates at $\nu_0=r$ by $\ell_r$, we have $\nu_0\ge r+1$ after stage $r$, and
$\nu_0=k$ after stage $k-1$.

\emph{Step 3: the output $\Psi(p)=\Opt(w)$.}  By Lemma~\ref{lem:duality} and Theorem~\ref{thm:A},
at $O(ms+m^{2})$ further gates the circuit outputs
\[
  \Psi(p)=\sum_{b=1}^{k}p_b+\beta_M(v)=\sum_{b=1}^{k}p_b-h_M(-v)=\Opt(w).
\]

\emph{Step 4: the size is $O(k^{3}m^{4}(s+m))$, and the family is uniform.}  The number of updates is at most
\[
  \begin{aligned}
  \sum_{r=1}^{k-1}\ell_r&\le(k-1)\bigl(1+k(2m-k-1)\bigr)\le(k-1)(1+2km)\\
  &=2k^{2}m-(2km-k+1)\le2k^{2}m.
  \end{aligned}
\]
An update of stage $r$ has $O(rm^{3}(s+m))$ gates by Lemma~\ref{lem:oneupdate},
and $\ell_r\le1+k(2m-k-1)\le2km$, so
\[
  \sum_{r=1}^{k-1}r\,\ell_r\le2km\sum_{r=1}^{k-1}r=k^{2}(k-1)m\le k^{3}m.
\]
So the size is $O(k^{3}m^{4}(s+m))$.  If every connected component of $M$ has
rank at most one, Corollary~\ref{cor:rankone} allows $\ell_r=r$ in Step 1, and then the size is $O(k^{3}m^{3}(s+m))$, as
\[
  \sum_{r=1}^{k-1}r\,\ell_r=\sum_{r=1}^{k-1}r^{2}\le k^{3}.
\]  The circuit depends only on $M$, the
blocks and the rank circuit.  Besides them it uses the rank $k=r_M(E)$ and the
connected components of $M$, which a polynomial-time algorithm obtains from the
rank circuit: $k$ by one evaluation, a base $B_0$ by the greedy algorithm, each
fundamental circuit $C(B_0,e)$ by $k$ evaluations, and the components as the
classes of the transitive closure of ``$e$ and $f$ lie in a common fundamental
circuit of $B_0$''~\cite{oxley11}.  So the family is uniform.
\end{proof}

\section{Proofs for Section~\texorpdfstring{\ref{sec:two}}{5}}
\label{app:two}

\begin{proof}[Proof of Lemma~\ref{lem:doubling}]
\emph{Step 1: $X\times\{1\}\cup(E\setminus X)\times\{2\}\in\Bases(N)\iff X\in\Bases(M_1)\cap\Bases(M_2)$.}  Since the bases of $M_2^{*}$ are the complements of
the bases of $M_2$,
\[
  \begin{aligned}
  X\times\{1\}\cup(E\setminus X)\times\{2\}\in\Bases(N)
  &\iff X\in\Bases(M_1),\ E\setminus X\in\Bases(M_2^{*})\\
  &\iff X\in\Bases(M_1)\cap\Bases(M_2),
  \end{aligned}
\]
and only the first copies carry weight, so the weight is $w(X)$.

\emph{Step 2: every rainbow base of $N$ has this form.}  $N$ has rank $\rho+(m-\rho)=m$, so a rainbow
base meets each of the $m$ blocks once, and it is the set of Step 1 for the set
$X$ of elements whose first copy it contains.
\end{proof}

\begin{proof}[Proof of Theorem~\ref{thm:D}]
\emph{Step 1: $\Opt(w)$ for $N$ is the minimum weight of a common base.}  If $M_1$ and $M_2$ have a common base, both have
rank $\rho$, and $N$ of Lemma~\ref{lem:doubling} has rank $m$ on $2m$ elements
with $m$ blocks, its rainbow bases corresponding to the common bases with equal
weights.

\emph{Step 2: a rank circuit of $N$ with $s_1+s_2+O(m)$ gates.}  With $r_i$ the rank function of $M_i$, for
$X_1,X_2\subseteq E$,
\[
  r_N\bigl(X_1\times\{1\}\cup X_2\times\{2\}\bigr)=r_1(X_1)+r_2^{*}(X_2)=r_1(X_1)+|X_2|+r_2(E\setminus X_2)-r_2(E),
\]
computed from the two rank circuits, the second on the negated indicator vector
of $X_2$, a population count, two additions and the subtraction of the constant
$r_2(E)$ on $O(\log m)$ bits, at $s_1+s_2+O(m)$ gates.

\emph{Step 3: the minimum weight of a common base with $O(m^{7}(s_1+s_2+m))$ gates.}  Theorem~\ref{thm:C} with $2m$ elements, $k=m$
blocks and the weights of Lemma~\ref{lem:doubling} gives the minimum weight of a
common base with
\[
  O\bigl(m^{3}(2m)^{4}(s_1+s_2+m)\bigr)=O\bigl(m^{7}(s_1+s_2+m)\bigr)
\]
gates.  The rank circuit of Step 2 is produced from the two given ones in
polynomial time, with $r_2(E)$ found by one evaluation, so the family is uniform
by Theorem~\ref{thm:C}.

\emph{Step 4: $X\cup Y\in\Bases(M_1')\cap\Bases(M_2')\iff X\in\mathcal I,\ |X|+|Y|=m$.}  Let $\mathcal I$ be the family of common
independent sets of $M_1$ and $M_2$, let $Z$ be a set of $m$ new elements, and
for $i=1,2$ let $M_i'$ be the matroid on $E\cup Z$ whose independent sets are the
sets $X\cup Y$ with $X$ independent in $M_i$, $Y\subseteq Z$ and
$|X|+|Y|\le m$: the truncation at $m$ of the direct sum of $M_i$ and the free
matroid on $Z$.  It has rank $m$ and the rank function
\[
  r_i'(S)=\min\bigl(r_i(S\cap E)+|S\cap Z|,\,m\bigr),
\]
which the rank circuit of $M_i$, a population count, an addition and a
comparison with $m$ compute at $s_i+O(m)$ gates.  Its bases are its independent
sets of size $m$, so for $X\subseteq E$ and $Y\subseteq Z$
\[
  X\cup Y\in\Bases(M_1')\cap\Bases(M_2')\iff X\in\mathcal I,\ |X|+|Y|=m,
\]
and $Z$ is a common base of $M_1'$ and $M_2'$.

\emph{Step 5: the maximum weight of a common independent set with $O(m^{7}(s_1+s_2+m))$ gates.}  Let $w'$ be $-w$ on $E$ and $0$ on $Z$, so
that $w'(X\cup Y)=-w(X)$.  Every $X\in\mathcal I$ has $|X|\le m=|Z|$, so by Step 4
\[
  \max_{X\in\mathcal I}w(X)=-\min_{X\in\mathcal I}\bigl(-w(X)\bigr)
  =-\min_{B\in\Bases(M_1')\cap\Bases(M_2')}w'(B).
\]
Step 3 for $M_1'$ and $M_2'$ on $2m$ elements, $m$ negations for $-w$ and one
for the output give this value with $O(m^{7}(s_1+s_2+m))$ gates, and the rank
circuits of $M_1'$ and $M_2'$ are produced in polynomial time as well.
\end{proof}

The circuit of Lemma~\ref{lem:augment} computes the largest size of a rainbow
independent set, that is, it solves unweighted matroid intersection with a
partition matroid.  Transversal matroids have rank functions of this kind, and
gammoids are restrictions of their duals, so Theorem~\ref{thm:A} applies to
both.

\begin{corollary}
\label{cor:transversal}
The following matroids $M$ on $m$ elements have Boolean rank circuits and $(\max,+,-)$-circuits for $h_M$ of polynomial size.
\begin{enumerate}
  \item Let $A_1,\ldots,A_n\subseteq E$, $|E|=m$, and $p=|A_1|+\cdots+|A_n|\le mn$.
        The partial transversals, the sets $X\subseteq E$ with an injective map
        $\phi:X\to[n]$ such that $x\in A_{\phi(x)}$ for every $x\in X$, are the
        independent sets of a matroid $M$, the transversal matroid.  It has a
        Boolean rank circuit of size at most $36(n+1)(p+1)^{2}(m+p+1)$, and $h_M$
        has a $(\max,+,-)$-circuit of size at most
        $120(m+1)(n+1)(p+1)^{2}(m+p+1)\le120(m+1)^{4}(n+1)^{4}$.
  \item Let $D$ be a directed graph on a set $V$ of $v$ vertices whose $a$ arcs are
        ordered pairs of vertices, and let $S,E\subseteq V$ with $|E|=m$.  A set
        $X\subseteq E$ is \emph{linked} from $S$ if there are directed paths, one
        from a vertex of $S$ to each $x\in X$, no two of which share a vertex; a
        vertex of $S\cap X$ is linked by the path without arcs.  The sets linked
        from $S$ are the independent sets of a matroid $M$, the gammoid.  It has a
        Boolean rank circuit of size at most $48(v+1)(v+a+1)^{2}(2v+a+1)$, and
        $h_M$ has a $(\max,+,-)$-circuit of size at most
        $160(m+1)(v+1)(v+a+1)^{2}(2v+a+1)\le160(v+1)^{8}$.
\end{enumerate}
\end{corollary}

\begin{proof}
\emph{Step 1: $r(U)=\nu(N_U)$.}  For $U\subseteq E$ let $N_U$ be the matroid on the
$p$ incidences $(e,i)$, $e\in A_i$, in which a set is independent when it contains
at most one incidence of each element and only incidences of elements of $U$, with
the blocks $R_i=\{(e,i):e\in A_i\}$.  Rainbow independent sets of $N_U$ are
matchings of elements of $U$ into distinct sets, and the partial transversals are
the independent sets of a matroid, as Edmonds and Fulkerson
showed~\cite{oxley11}, so
\[
  r(U)=\max\{|X|:X\subseteq U\text{ a partial transversal}\}=\nu(N_U).
\]

\emph{Step 2: the bounds of item 1.}  The set $J$ is independent in
$N_U$ if and only if
\[
  |J|=r_{N_U}(J)=\bigl|\{e\in U:(e,i)\in J\text{ for some }i\}\bigr|,
\]
a test of at most $28p+23m+1$ gates with the bits of $U$ as extra inputs: one
disjunction per incidence, one conjunction per element and two population
counts.  Lemma~\ref{lem:augment} on the $p$ incidences and $n$ blocks and a
population count give $r(U)$ within $36(n+1)(p+1)^{2}(m+p+1)$ gates.
Theorem~\ref{thm:A} with $k\le m$ gives the second bound, and $p\le mn$ gives
$p+1,\ m+p+1\le(m+1)(n+1)$ and the last one.

\emph{Step 3: $X$ is linked from $S$ $\iff$ $X$ is independent in $T^{*}$.}  Let $T$ be
the transversal matroid on $V$ of $A_y=\{y\}\cup\{z:z\to y\}$, $y\in V\setminus S$.
By a theorem of Ingleton and Piff~\cite{oxley11},
\[
  \begin{aligned}
  X\text{ is linked from }S&\iff X\cap B=\varnothing\text{ for some }B\in\Bases(T)\\
  &\iff X\text{ is independent in }T^{*}.
  \end{aligned}
\]
So the gammoid is the restriction of $T^{*}$ to $E$.

\emph{Step 4: the bounds of item 2.}  The family $(A_y)$ has at most $v$ sets
and $v+a$ incidences, so Step 2 for $T$, the dual rank circuit of the proof of Corollary~\ref{cor:classes}(4) on $V$, and the input bits placed on $E$ ($0$ on
$V\setminus E$) give a rank circuit with
\[
  \begin{aligned}
  &36(v+1)(v+a+1)^{2}(2v+a+1)+(v+1)+11(v+1)\lceil\log_2(v+1)\rceil\\
  &\qquad\le48(v+1)(v+a+1)^{2}(2v+a+1)
  \end{aligned}
\]
gates.  Theorem~\ref{thm:A} with $k\le m\le v$ and $a\le v^{2}$ gives the other
two bounds.
\end{proof}

\section{Formal Verification}
\label{app:checks}

A Lean~4 development with Mathlib, both at version 4.32.0, proves the results of
Table~\ref{tab:lean} for the matroids of Mathlib's \texttt{Matroid} on
$\{0,\ldots,m-1\}$, with explicit bounds in place of $O(\cdot)$.  The main theorems
depend on no axioms beyond \texttt{propext}, \texttt{Classical.choice} and
\texttt{Quot.sound}, and the build reports no warnings.

\begin{table}[ht]
\caption{The results proved in Lean.}
\label{tab:lean}
\centering
\renewcommand{\arraystretch}{1.25}
\begin{tabular}{@{}>{\raggedright\arraybackslash}p{0.43\textwidth}
                >{\raggedright\arraybackslash}p{0.53\textwidth}@{}}
\toprule
Results & Lean statement\\
\midrule
Definition~\ref{def:ti} & the two expressions agree\\
Theorems~\ref{thm:compiler}, \ref{thm:A}, \ref{thm:B}; Lemmas~\ref{lem:greedy},
\ref{lem:minbases}, equation~\eqref{eq:rank-mu}; Corollaries~\ref{cor:levels},
\ref{cor:classes}, \ref{cor:integer} & as stated\\
Corollary~\ref{cor:count} & for $m\ge128$, at most a fraction
$2^{-\lfloor2^{m}/m^{3}\rfloor}$ of the matroids on $m$ labeled elements have a
$(\max,+,-)$-circuit of size less than $2^{m/2}/m^{2}$ that computes $h_M$ on
$\{0,1\}^{E}$\\
Lemmas~\ref{lem:duality}, \ref{lem:updint}, \ref{lem:events},
\ref{lem:augment}--\ref{lem:doubling}, \ref{lem:path}--\ref{lem:crossing};
Propositions~\ref{prop:pos}, \ref{prop:obstacles}(1), \ref{prop:posl},
\ref{prop:star}; Theorems~\ref{thm:validity}, \ref{thm:potential},
\ref{thm:stagelength}; Corollaries~\ref{cor:rankone}, \ref{cor:transversal} & as
stated\\
Proposition~\ref{prop:obstacles}(2) & the reduction, Steps 2 to 4, for every
multigraph $H$ and every $h\ge|V(H)|,|Q|$\\
Theorem~\ref{thm:C} & size at most $3328\,k^{3}m^{4}(s+m)+2$, and
$1390\,k^{3}m^{3}(s+m)+2$ when every component of $M$ has rank at most one\\
Theorem~\ref{thm:D} & size at most $3328\cdot1200\,m^{7}(s_1+s_2+m)+2$ for common
bases and $3328\cdot653313\,m^{7}(s_1+s_2+m)+3$ for common independent sets\\
\midrule
Not formalized & uniformity; the complexity classes in
Proposition~\ref{prop:obstacles}(2), where the NP-completeness of rainbow path is
cited\\
\bottomrule
\end{tabular}
\end{table}


\end{document}
